% concatenated from Arxiv_v2.tex
%%%%%%%% ICML 2026 EXAMPLE LATEX SUBMISSION FILE %%%%%%%%%%%%%%%%%

\documentclass{article}

\input{mystyle.sty}



% Recommended, but optional, packages for figures and better typesetting:
\usepackage{microtype}
\usepackage{graphicx}
\usepackage{subcaption}
\usepackage{booktabs} % for professional tables

\usepackage{subcaption}  % in preamble


% hyperref makes hyperlinks in the resulting PDF.
% If your build breaks (sometimes temporarily if a hyperlink spans a page)
% please comment out the following usepackage line and replace
% \usepackage{icml2026} with \usepackage[nohyperref]{icml2026} above.
% \usepackage[dvipsnames]{xcolor}        % For defining colors
% \usepackage{hyperref}
% \hypersetup{
%     colorlinks=true,       % Enable colored links
%     linkcolor=blue,       % Color for internal links
%     urlcolor=blue,    % Color for external links (like URLs)
% }

% Attempt to make hyperref and algorithmic work together better:
\newcommand{\theHalgorithm}{\arabic{algorithm}}
\setlength{\textfloatsep}{1pt plus 1.0pt minus 2.0pt}
\setlength{\floatsep}{1pt plus 1.0pt minus 2.0pt}
\setlength{\intextsep}{1pt plus 1.0pt minus 2.0pt}

% Use the following line for the initial blind version submitted for review:
% \usepackage{icml2026}

\usepackage[dvipsnames]{xcolor}        % For defining colors
\usepackage{hyperref}
% \hypersetup{
%     colorlinks=true,       % Enable colored links
%     linkcolor=blue,       % Color for internal links
%     urlcolor=blue,    % Color for external links (like URLs)
%   citecolor=blue,
% }
\hypersetup{
  colorlinks=true,
  linkcolor=MidnightBlue,
  citecolor=MidnightBlue,
  urlcolor=MidnightBlue
}

% For preprint, use
\usepackage[preprint]{icml2026}

% If accepted, instead use the following line for the camera-ready submission:
% \usepackage[accepted]{icml2026}

\usepackage{amsmath}
\usepackage{amssymb}
\usepackage{mathtools}
\usepackage{amsthm}
\usepackage{dsfont}

\newcommand{\defi}[1]{{\bf #1}}

% if you use cleveref..
\usepackage[capitalize,noabbrev]{cleveref}

\usepackage{mathtools}
\mathtoolsset{showonlyrefs}
\usepackage{enumitem}

% % 2. 'algpseudocodex' replaces 'algorithmic'. 
% %    - [noEnd=true] hides the "end for" text.
% %    - It draws vertical lines by default.
% \usepackage[noEnd=true]{algpseudocodex}


%%%%%%%%%%%%%%%%%%%%%%%%%%%%%%%%
% THEOREMS
%%%%%%%%%%%%%%%%%%%%%%%%%%%%%%%%
\theoremstyle{plain}
\newtheorem{theorem}{Theorem}[section]
\newtheorem{proposition}[theorem]{Proposition}
\newtheorem{lemma}[theorem]{Lemma}
\newtheorem{corollary}[theorem]{Corollary}
% \theoremstyle{definition}
\newtheorem{definition}[theorem]{Definition}
\newtheorem{assumption}[theorem]{Assumption}
% \theoremstyle{remark}
\newtheorem{remark}[theorem]{Remark}

% Todonotes is useful during development; simply uncomment the next line
%    and comment out the line below the next line to turn off comments
%\usepackage[disable,textsize=tiny]{todonotes}
% \usepackage[textsize=tiny]{todonotes}




% EXTRA MACROS FOR RECONSTRUCTION PART
\newcommand{\muNext}{\mu'}
\newcommand{\hZ}{\hat{\mathbf{z}}}
\DeclareMathOperator*{\argmin}{arg\,min}

% The \icmltitle you define below is probably too long as a header.
% Therefore, a short form for the running title is supplied here:
\icmltitlerunning{Population-Aware IL in MFGs with Common Noise}

\begin{document}

\twocolumn[
  \icmltitle{Population-Aware Imitation Learning in Mean-field Games with Common Noise}

  % It is OKAY to include author information, even for blind submissions: the
  % style file will automatically remove it for you unless you've provided
  % the [accepted] option to the icml2026 package.

  % List of affiliations: The first argument should be a (short) identifier you
  % will use later to specify author affiliations Academic affiliations
  % should list Department, University, City, Region, Country Industry
  % affiliations should list Company, City, Region, Country

  % You can specify symbols, otherwise they are numbered in order. Ideally, you
  % should not use this facility. Affiliations will be numbered in order of
  % appearance and this is the preferred way.
  \icmlsetsymbol{equal}{*}

  \begin{icmlauthorlist}
    \icmlauthor{Gr\'egoire Lambrecht}{cds}
    \icmlauthor{Mathieu Lauri\`ere}{ecnu}
  \end{icmlauthorlist}

  \icmlaffiliation{cds}{Center for Data Science, New York University and NYU Shanghai, New York, United States of America}
  \icmlaffiliation{ecnu}{Shanghai Center for Data Science; NYU-ECNU Institute of Mathematical Sciences at NYU Shanghai; NYU Shanghai, Shanghai, People’s Republic of China}

  \icmlcorrespondingauthor{Gr\'egoire Lambrecht}{gl3048@nyu.edu}
  \icmlcorrespondingauthor{Mathieu Lauri\`ere}{ mathieu.lauriere@nyu.edu}

  % You may provide any keywords that you find helpful for describing your
  % paper; these are used to populate the "keywords" metadata in the PDF but
  % will not be shown in the document
  \icmlkeywords{Machine Learning, ICML}

  \vskip 0.3in
]

% this must go after the closing bracket ] following \twocolumn[ ...

% This command actually creates the footnote in the first column listing the
% affiliations and the copyright notice. The command takes one argument, which
% is text to display at the start of the footnote. The \icmlEqualContribution
% command is standard text for equal contribution. Remove it (just {}) if you
% do not need this facility.

% Use ONE of the following lines. DO NOT remove the command.
% If you have no special notice, KEEP empty braces:
\printAffiliationsAndNotice{}  % no special notice (required even if empty)
% Or, if applicable, use the standard equal contribution text:
% \printAffiliationsAndNotice{\icmlEqualContribution}

\begin{abstract}
% We consider two imitation learning problems for Mean Field Games (MFGs) with common noise. MFGs are a powerful framework to study equilibria of systems composed of a large number of heterogeneous agents. Imitation learning for MFGs generally aims to give guarantees on the recovery of a Nash equilibrium when one can reproduce an equilibrium policy. Existing results show that Behavioral Cloning (BC) and Adversarial Learning (ADV) are two relevant proxies for these problems.

% Beyond recovering a policy, we also address the problem of performance against an expert population. This corresponds to the situation where one tries to learn an equilibrium policy and then plays against the expert population. In the presence of common noise, the density generated by the agents during the game becomes random. As a result, agents need to consider richer policies that depend on the population.

% We show that in this context, BC and ADV are still powerful. We illustrate our results numerically by showing that, in many situations, if one tries to learn an equilibrium by regressing it with a population-independent policy, the result is far from an equilibrium and performs poorly against experts. Since in many real-world situations it is difficult to determine whether common noise is present or not, our results show that it is better to consider population-dependent policies for imitation learning in MFG in order to avoid being fooled by randomness coming from common noise.


Mean Field Games (MFGs) provide a powerful framework for modeling the collective behavior of large populations of interacting agents. In this paper, we address the problem of Imitation Learning (IL) in MFGs subject to \emph{common noise}, where the population distribution evolves stochastically. This stochasticity compels agents to adopt \emph{population-aware policies} to respond to aggregate shocks. We formulate two distinct learning objectives: recovering a Nash equilibrium and maximizing performance against an expert population. We investigate two imitation proxies: Behavioral Cloning (BC) and Adversarial (ADV) divergence. We then establish finite-sample error bounds showing that minimizing these proxies effectively controls both the policy's exploitability and its performance gap relative to the expert. 
% Our theoretical analysis encompasses general dynamics ($L_P \geq 0$). 
Furthermore, we propose a numerical framework using generalized Fictitious Play and Deep Learning to compute expert population-aware policies. Through experiments on three environments %—a stylized 2-state model, a ``Beach Bar'' congestion problem, and a ``Night Club'' social choice scenario—
we demonstrate that standard population-unaware policies fail to capture the equilibrium dynamics. Our results highlight that learning population-aware policies is crucial to avoid being misled by the randomness inherent in common noise.

% \gl{"Fooled by randomness" is famous? book by Nassim Nicholas Taleb.} \ml{I didn't know about it but let me know if you want to stick to it}
% \gl{The abstract works to me ! May be, I should wait for when I will have enough citations before making any strange easter egg haha.}

\end{abstract}

% \ml{We should harmonize terminology: population-aware, population-dependent, adaptive, Master policies.} \gl{Switching to population-aware}

\section{Introduction}

% \ml{TODO}

Mean Field Games (MFGs), introduced independently by~\citep{MR2295621} and~\citep{MR2346927}, have emerged as a fundamental framework for analyzing systems with a very large number of rational, indistinguishable agents. By taking the limit as the number of agents $N \to \infty$, the interactions are simplified into a game between a single representative agent and the aggregate population distribution. 
While classical MFGs assume that the population distribution evolves deterministically (see e.g.~\citep{MR3134900,MR3752669}), many real-world applications are subject to systemic shocks affecting the entire population simultaneously. Examples include macro-economics~\cite{lasry2008application,gueant2010mean,vu2025heterogenous}, financial markets \cite{carmona2015mean,carmona2017mean,firoozi2018mean,fu2021mean,burzoni2023mean,bassou2024mean}, price formation~\cite{graber2016linear,gomes2023machine},  electricity management~\cite{alasseur2020extended,lavigne2023decarbonization,escribe2024mean,dumitrescu2024energy} and crowd motion~\cite{achdou2018mean,perrin2020fictitious}. These scenarios are modeled by MFGs with \defi{common noise}, where the mean field becomes a stochastic process adapted to the common noise filtration, see~\citep{carmona2016mean,MR3753660} for the theoretical background. Solving such games is significantly more challenging, because the mean field is \emph{stochastic}. 

Although the theoretical framework is developed and various applications have been proposed, so far there are no works about learning common noise MFG models from data. This work is a first step in this direction. Here, we focus on the problem of \defi{Imitation Learning (IL)} in MFGs with common noise. We consider a scenario where a learner observes a population of experts playing a Nash equilibrium but does not know the reward function driving their behavior. The goal is to learn a policy that imitates the expert, ideally recovering a Nash equilibrium and matching the expert's performance.

\noindent\textbf{Related Work:} 
IL and Inverse Reinforcement Learning (IRL) are well-established in single-agent settings~\citep{ng2000algorithms,ho2016generative}. In the context of Multi-Agent Reinforcement Learning (MARL) and MFGs, recent works have begun to address the learnability of equilibria. To the best of our knowledge, \citet{subramanian2019reinforcement},  \citet{guo2019learningMFG}, and \citet{elie2020convergence} were the first to develop and analyze RL for MFGs, aiming at computing Nash equilibria using RL methods and a simulator for the MFG model. RL for MFGs with common noise was studied in~\citep{perrin2020fictitious,wu2025population}. Here, we consider the \emph{converse} problem: learning a policy by looking at data generated by an expert. Along this line, \citet{ramponi2023imitation} provide a theoretical foundation for IL in \emph{deterministic} MFGs. We generalize these results to the case with common noise.\citep{Anahtarci2025MaxCausalEntropyMFG,Chen2022IndividualIRLMFG,Zhao2025MeanFieldCorrelatedIL} proposed methods to learn policies and \cite{Chen2024MetaIRLMFG} methods to learn rewards, but so far there are no such works with \emph{common noise}.

\noindent\textbf{Challenges:} 
First, common noise breaks the deterministic nature of the mean-field flow, requiring a \textbf{probabilistic} analysis of the error. Second, this implies that vanilla policies functions of the individual state only are sub-optimal, 
requiring agents to adopt \textbf{adaptive} policies that react to the fluctuating distribution of the population.

\noindent\textbf{Contributions:} 
We provide the first comprehensive \textbf{theoretical} and \textbf{algorithmic} framework for imitation learning in MFGs with \textbf{common noise}. 
Our main contributions are:
\begin{enumerate}%[noitemsep,topsep=0pt]
    \item \textbf{Theoretical Guarantees:} We analyze two imitation metrics: a local Behavioral Cloning (BC) proxy and a global Adversarial (ADV) proxy. We prove that minimizing these proxies guarantees the learned policy is an approximate Nash equilibrium and achieves a reward comparable to the expert. %Crucially, our bounds account for the Lipschitz continuity of the environment dynamics ($L_P$) and the expert policy ($L_E$) with respect to the stochastic mean field.
    \item \textbf{Algorithmic Framework for Equilibria:} Computing a ``ground truth'' expert policy in common noise settings is non-trivial. We generalize the Fictitious Play algorithm to the common noise setting. To resolve the issue that Fictitious Play yields a sequence of policies rather than a single deployable one, we introduce a method to distill this sequence into a single population-dependent policy.
    \item \textbf{Population-aware Imitation:} We highlight the distinction between \textit{vanilla} policies (ignoring the mean field) and \textit{adaptive} policies (observing it). We propose an interactive imitation learning algorithm that trains agents to internalize the relationship between common noise shocks, population shifts, and optimal actions.
\end{enumerate}

The remainder of the paper is organized as follows. Sec.~\ref{sec:pb-setting} formalizes the framework. Sec.~\ref{sec:il-results} presents our main theoretical bounds linking imitation proxies to equilibrium quality. Sec.~\ref{sec:numerics} details the numerical methods and experimental results. Sec.~\ref{sec:ccl} concludes the paper.

% \vspace{-10pt}

% \gl{Same than for the abstract, it works for me ! May be we should remve the '--'.}

\section{Problem Setting} 
\label{sec:pb-setting}

\subsection{Mean Field Dynamics and Nash Equilibria} 

\noindent\textbf{Notation.} $\Delta_{\cS}$ denotes the set of probability measures over any finite set $\cS$. For any integer $N$, $[N] = \{0, \dots, N\}$ and $[N]^* = [N] \setminus \{0\}$.

\noindent\textbf{Main components.} 
A \defi{Mean Field Game (MFG)} with common noise is denoted by a tuple $\cM = (\cX, \cA, \cE^0, \nu^0, P, r, H, \rho_0)$, where $\cX$ is a finite state space, $\cA$ is a finite action space, $\cE^0$ is a Polish space for the \defi{common noise} impacting the whole population, $\nu^0$ is a probability measure for the common noise, while $P: \cX \times \cA \times \Delta_{\cX} \times \cE^0 \to \Delta_{\cX}$ is a stochastic kernel for the agents' dynamics, the mapping $r: \cX \times \cA \times \Delta_{\cX} \to \mathbb{R}$ is a measurable reward function, $H \in \mathbb{N}^*$ is the time horizon, and $\rho_0 \in \Delta_{\cX}$ is the initial probability measure. We define a \defi{strategy} as a measurable mapping from $\cX \times \Delta_{\cX}$ to $\Delta_{\cA}$, and a \defi{policy} as a sequence of $H$ strategies. A strategy is denoted by $\pi_0$ and the set of all such strategies is denoted by $\Pi_0$. A policy is denoted by $\pi = (\pi_0, \dots, \pi_{H-1})$, and the set of all such policies is denoted by $\Pi$. Elements of $\Delta_{\cX}$ and $(\Delta_{\cX})^H$ are referred to as \defi{mean field} and \defi{mean field sequences} respectively.

\noindent\textbf{Distribution evolution.} We consider \emph{two types} of distributions: the distribution of the whole population using a given policy, and the distribution of a single agent's state using a given policy when the population is possibly using a different policy. 
For any pair of population distributions $\rho, \rho' \in \Delta_\cX$, any strategy $\pi_0 \in \Pi_0$, we define the mean-field transition operator $\Phi(\rho', \pi_0; \rho, e^0)$ as
% \vspace{-9pt}
$$
    \Phi(\rho', \pi_0; \rho, e^0) = \sum_{x \in \cX, a \in \cA} \rho'(x) \pi_0 (a \vert x, \rho) P(\cdot \vert x, a, \rho, e^0).
    % \vspace{-10pt}
$$
This operator $\Phi$ represents the evolution of the probability distribution of an agent, given that the agent's state distribution was $\rho'$, the agent used strategy $\pi_0$, the total population was distributed according to $\rho$, and the common noise value is $e^0$.
We consider an independent and identically distributed (i.i.d.) sequence of random variables $\varepsilon^0_1, \dots, \varepsilon^0_{H-1} \in \cE^0$ for the common noise. We denote the common noise trajectory up to time $H-1$ as $\varepsilon^0 = (\varepsilon^0_1, \dots, \varepsilon^0_{H-1})$. For any policy $\pi \in \Pi$, we define the mean field sequence $\rho^{\pi, \varepsilon^0} = (\rho^{\pi, \varepsilon^0}_0, \dots, \rho^{\pi, \varepsilon^0}_{H-1})$ by the recursion: $\rho^{\pi, \varepsilon^0}_0 = \rho_0$ and for $n \in [H-2]$,
% \vspace{-5pt}
\begin{equation*}
        \rho^{\pi, \varepsilon^0}_{n+1} = \Phi(\rho^{\pi, \varepsilon^0}_n, \pi_n \,; \, \rho^{\pi, \varepsilon^0}_n, \varepsilon^0_{n+1}).
        % \vspace{-7pt}
\end{equation*}
This sequence represents the evolution of the population distribution when all agents follow policy $\pi$ and are affected by the common noise trajectory $\varepsilon^0$. For any pair of policies $(\pi, \pi') \in \Pi^2$, we define the mean field sequence $\rho^{\pi, \pi', \varepsilon^0} = (\rho^{\pi, \pi', \varepsilon^0}_0, \dots, \rho^{\pi, \pi', \varepsilon^0}_{H-1})$ by: $\rho^{\pi, \pi', \varepsilon^0}_0 = \rho_0$ and for $n \in [H-2]$, 
% \vspace{-6pt}
\begin{equation*}
        \rho^{\pi, \pi', \varepsilon^0}_{n+1} = \Phi(\rho^{\pi, \pi'_n, \varepsilon^0}_{n}, \pi'; \rho^{\pi, \varepsilon^0}_n, \varepsilon^0_{n+1}).
% \vspace{-7pt}
\end{equation*}
This sequence represents the evolution of the state distribution of a \emph{single agent} following policy $\pi'$ when all other agents follow policy $\pi$ and are affected by the common noise trajectory $\varepsilon^0$. %In the notation for mean-field trajectories, we will often omit 

\noindent\textbf{Best response and Nash equilibrium.} 
When the population of agents follows the policy $\pi$, a single agent wishes to find a \textbf{best response} to $\pi'$, i.e., a policy maximizing the \textbf{value function}:
% \vspace{-13pt}
$$
    V(\pi', \pi) = \mathbb{E} \left[ \sum_{n = 0}^{H-1} r(X_n, \alpha_n, \rho^{\pi, \varepsilon^0}_n) \right],
% \vspace{-7pt}
$$
where the process is defined by $X_0 \sim \rho_0$, $X_{n+1} \sim P(\cdot \vert X_n, \alpha_n, \rho^{\pi, \varepsilon^0}_n, \varepsilon^0_{n+1})$, and $\alpha_n \sim \pi'(\cdot \vert X_n, \rho^{\pi ,\varepsilon^0}_n)$. Note that under these conditions, the law of the state $X_n$ (conditioned on the common noise) is $\rho^{\pi, \pi', \varepsilon^0}_n$. We denote the law of state-action pair $(X_n, \alpha_n)$ by $\mu^{\pi, \pi', \varepsilon^0}_n \in \Delta_{\cX\times \cA}$.
If $\pi = \pi'$, we will simply denote $\mu^{\pi, \varepsilon^0}_n = \mu^{\pi, \pi', \varepsilon^0}_n$.  In the sequel, we will often omit to write explicitly the common noise $\varepsilon^0$, but the sequences $\rho$ and $\mu$ are random variables due to $\varepsilon^0$.


\begin{definition}
    For $\pi \in \Pi$, the set of best responses is defined by $\text{BR}(\pi) = \arg \max_{\pi' \in \Pi} V(\pi', \pi)$.
\end{definition}
We will use the concept of Nash equilibrium: a situation where no agent has an incentive to unilaterally deviate from their chosen policy.  %In our mean-field setting, a Nash equilibrium corresponds to a condition of mutual best response between the population and the individual agent.
\begin{definition}
    A policy $\pi \in \Pi$ is a \defi{(mean field) Nash equilibrium} if $\pi \in \text{BR}(\pi)$.
\end{definition}
This is a fixed point-condition: it implies that when the entire population follows policy $\pi$, the resulting mean-field flow $\rho^{\pi}$ creates a reward and transition structure such that the same policy $\pi$ remains optimal for any individual agent.


\subsection{Imitation in Unknown Reward Environments}
Our objective is as follows. Assume that an observer sees realizations of an expert policy $\pi^E$ that constitutes a Nash equilibrium for a MFG $\cM$. We further assume that the instantaneous reward function $r$ is unknown to the observer. The observer tries to imitate $\pi^E$ and computes a candidate policy $\pi^A$. We wish to study two questions:
\begin{enumerate}[
  label=\textbf{(Q\arabic*)},
  ref=\textbf{(Q\arabic*)}]
%   noitemsep,
%   topsep=0pt
% ]
    \item\label{theory-q1} \textit{How far is $\pi^A$ from being a Nash equilibrium?}
    \item\label{theory-q2} \textit{How far is $\pi^A$ from being a best response to $\pi^E$?}
\end{enumerate}
The first corresponds to a scenario where a whole population would be using $\pi^A$ while the second corresponds to a scenario where the population is playing $\pi^E$ and a single agent uses policy $\pi^A$.
% quantify how far a candidate policy $\pi^A$ is from satisfying two properties:
% \begin{enumerate}
%     \item Being a Nash equilibrium itself,
%     \item Performing as well as $\pi^E$ in a scenario where the rest of the population follows the expert policy $\pi^E$.
% \end{enumerate}
Specifically, we aim to answer these questions in quantitative way using two divergences calculated relative to the expert policy $\pi^E$. The first, $\delta^{BC}$, captures the local discrepancy in decision-making (Behavioral Cloning), while the second, $\delta^{ADV}$, captures the global divergence in the induced state-action distributions (Adversarial).

We define the \defi{Behavioral Cloning (BC) divergence} at time $n$ as:
% \vspace{-5pt}
\begin{equation*}
    \delta^{BC}_n = \mathbb{E} \left[ \EE_{X \sim \rho_n^{\pi^E}}\left( \| \pi^A_n( X, \rho_n^{\pi^E}) - \pi^E_n(X, \rho_n^{\pi^E})\|_1  \right)\right],
    % \vspace{-6pt}
\end{equation*}
where the expectation $\mathbb{E}$ is taken over $\varepsilon^0_{1:H-1}$. 
We define the \defi{Adversarial (ADV) proxy} at time $n$ as:
% \vspace{-6pt}
\begin{equation*}
    \delta^{ADV}_n = \mathbb{E} \left[ \| \mu^{\pi^E}_n - \mu^{\pi^A}_n \|_1 \right].
    % \vspace{-8pt}
\end{equation*}
We denote the maximum divergences over the horizon as $\delta^{ADV} = \max_{n \in [H-1]} \delta^{ADV}_n$ and $\delta^{BC} = \max_{n \in [H-1]} \delta^{BC}_n$. We will often use the notation $\rho^E = \rho^{\pi^E}, \mu^E = \mu^{\pi^E}$ and $\rho^A = \rho^{\pi^A}, \mu^A = \mu^{\pi^A}$.

To answer the two questions above, we introduce the following regularity assumptions.

\begin{assumption}
    \label{assumption:Lipschitz}
    $r$ and $P$ are Lipschitz-continuous w.r.t. the mean-field with constant $L_r, L_P \ge 0$. 
\end{assumption}

\begin{assumption}
    \label{assumption:rmax}
    There exists $r_{\max} > 0$ satisfying for any $\pi \in \Pi$:
    % \vspace{-7pt}
    \begin{equation*}
        \PP(|r(x, a, \rho^{\pi, \varepsilon^0}_n)| \leq r_{\max}) = 1, \quad\forall (x, a) \in \cX \times \cA.
        % \vspace{-7pt}
    \end{equation*}
\end{assumption}

\begin{definition}[Mean-Field Lipschitz Policies]
    \label{definition:piE-Lipschitz}
    We define the class of \defi{Mean-Field Lipschitz Policies}, denoted by $\Pi^\cM_{\text{lip}}$, as the set of policies $\pi$ for which there exists a constant $L_E \ge 0$ such that for any sequence of policies $(\pi', \pi'', \pi^1, \pi^2, \pi^3, \pi^4) \in \Pi^6$ and any time step $n \in [H-1]$, the following inequality holds: 
    % \ml{the 6 other policies don't seem to play any role, we can just take 3 mean field sequences, isn't it?} \gl{We just need $\pi$ to be lipshitz on policies generated by the MFG. So the 6 policies play the role: $\rho_n^{\pi',\pi''}, \rho_n^{1,2}, \rho_n^{3,4'}$ are generated by the MFG. We could define $\{\rho^{\pi^1, \pi^2}_n(e^0_{1:n}), \; n\in [H-1], \; \pi^1, \pi^2\in \Pi, \; e^0_{k} \in supp(\nu^0)\}$ and say the propoerties for some rho in that set.}
    % \vspace{-9pt}
    \begin{align*}
    &\mathbb{E} \left[ 
        \EE_{X\sim\rho_n^{\pi',\pi''}} \left(
      \|\pi(x,\rho_n^{\pi^1,\pi^2})
        -\pi( x,\rho_n^{\pi^3,\pi^4})\|_1 \right)
    \right] \\
    &\qquad \le L_E \mathbb{E} \left[ 
        \|\rho_n^{\pi^1,\pi^2} - \rho_n^{\pi^3,\pi^4}\|_1 
    \right].
    % \vspace{-7pt}
    \end{align*}
\end{definition}
We will assume that our policy space $\Pi$ is restricted to a compact subset of $\Pi^\cM_{\text{lip}}$. This is a well-founded assumption, as the set of policies that are Lipschitz-continuous with a fixed constant $L_E$ with respect to the mean-field is indeed compact. From this point forward, we will write $\Pi$ to refer to this compact set of Lipschitz policies, utilizing an abuse of notation for the sake of brevity.
\begin{definition}
    The \defi{exploitability} of $\pi \in \Pi$ is defined as:
    $
        \cE(\pi) = \max_{\pi' \in \Pi} V(\pi', \pi) - V(\pi, \pi).
    $
\end{definition}
% \vspace{-6pt}
Exploitability serves as a direct measure of the proximity to be an equilibrium: $\pi$ is a Nash equilibrium iff $\mathcal{E}(\pi) = 0$, while a value close to zero indicates that the policy is nearly optimal for each agent. 
\begin{definition}
    For $\epsilon \ge 0$, a policy $\pi \in \Pi$ is an \defi{$\epsilon$-Nash equilibrium} if $\mathcal{E}(\pi) \leq \epsilon$.
\end{definition}
% \vspace{-6pt}
In the context of recovering an equilibrium from the policy $\pi^E$, we want to control $\cE(\pi) $ with $\delta^{BC}$ and $\delta^{ADV}$.
% \vspace{-10pt}





\section{Imitation Learning Results}
\label{sec:il-results}

We first address~\ref{theory-q1}. Due to space constraints, we provide the results and the main ideas of the proofs in Appendix~\ref{appendix:theory}.


\begin{theorem}
    \label{theorem:BC_LPge0}
    Suppose that Assm~\ref{assumption:Lipschitz} and~\ref{assumption:rmax} hold. We have
    \begin{equation*}
        \begin{split}
            \cE(\pi^A) &\leq \delta^{BC}\Big[r_{\max} H^2 
            \\
            &\qquad+ 2\frac{(1 + L_P + L_E)^H}{(L_P + L_E)^2} (L_r + r_{\max}(L_E + 1)\Big].
        \end{split}
    \end{equation*}
\end{theorem}

\begin{theorem}
    \label{theorem:ADV_LPge0}
    Suppose Assm~\ref{assumption:Lipschitz} and~\ref{assumption:rmax} hold. Then,
    \begin{equation}
        \begin{split}
            \cE(\pi^A) 
            &\leq \delta^{ADV}\Big[H(2L_r + 3L_E r_{\max} + r_{\max}) 
            \\
            &\qquad + 3r_{\max}H^2(L_E + L_P) \Big].
        \end{split}
    \end{equation}
\end{theorem}

\begin{remark}
    Our results are consistent with the ones established in \citep{ramponi2023imitation}. By letting $L_E = 0$ in Theorem~\ref{theorem:BC_LPge0}, one obtains \citep[Theorem~3]{ramponi2023imitation}. We recover \citep[Theorem~5]{ramponi2023imitation} when $L_E = 0$ in Theorem~\ref{theorem:ADV_LPge0}.
\end{remark}

\begin{remark}
    We do not address the case where $L_P = 0$ and $L_E = 0$. In this specific setting, the dynamics and the expert policy are independent of the population distribution. Under these conditions, it is straightforward to observe that the results from \citet[Theorem~1]{ramponi2023imitation} and \citet[Theorem~2]{ramponi2023imitation} directly apply.
\end{remark}

We then turn to~\ref{theory-q2} and give the guarantees of the performance of the policy $\pi^A$ versus $\pi^E$.
\begin{theorem}
    \label{theorem:BC_LPge0_reward}
    Suppose that Assm~\ref{assumption:rmax} holds. We have
    \begin{equation}
        \begin{split}
            \vert V(\pi^E, \pi^E) - V(\pi^A, \pi^E) \vert \leq H^2 \delta^{BC} r_{\max}.
        \end{split}
    \end{equation}
\end{theorem}

\begin{theorem}
    \label{theorem:ADV_LPge0_reward}
    Suppose Assm~\ref{assumption:Lipschitz} and~\ref{assumption:rmax} hold. We have
    \begin{equation}
        \begin{split}
            \vert V(\pi^E, \pi^E) - V(\pi^A, \pi^E) \vert 
            &\leq r_{\max}\delta^{ADV}\Big[H(L_E +1) 
            \\
            &\qquad + H^2 (L_P + L_E)\Big].
        \end{split}
    \end{equation}
\end{theorem}

% \vspace{-8pt}
\begin{remark}
    In Theorem~\ref{theorem:ADV_LPge0_reward}, Assumption~\ref{assumption:Lipschitz} is required only for the Lipschitz continuity of the transition kernel $P$, but not for the reward $r$. 
\end{remark}
% \vspace{-6pt}
\noindent\textbf{Results Interpretation and Analysis:}
Our results demonstrate that answering Question \ref{theory-q1} (Equilibrium Recovery) is significantly more subtle than answering Question \ref{theory-q2} (Performance Matching). 
% \ml{not clear what it means that a ``question'' is more robust than another question}

Robustness of Performance~\ref{theory-q2}: As shown in Theorem \ref{theorem:BC_LPge0_reward}, the performance gap is bounded by $\delta^{BC}$ without requiring Lipschitz continuity for the reward ($L_r$) or transition kernel ($L_P$). Even in complex or discontinuous environments, local action mimicry is sufficient to achieve expert-level utility.

Fragility of Equilibrium~\ref{theory-q1}: This stability does not hold for equilibrium recovery. An agent may achieve high rewards while remaining far from a Nash equilibrium. Theorem \ref{theorem:BC_LPge0} shows that exploitability $\cE(\pi^A)$ is sensitive to social reactivity ($L_E$) and environmental coupling ($L_P$), resulting in exponential error growth.

Finally, our results suggest that the ADV divergence is more robust for solving \ref{theory-q1} (Equilibrium Recovery) than BC divergence, because it transforms the error propagation from exponential (Theorem~\ref{theorem:BC_LPge0}) to polynomial (Theorem~\ref{theorem:ADV_LPge0}).




\section{Numerical Experiments}
\label{sec:numerics}

\subsection{Objectives and Method}
We aim to illustrate the control of the Behavioral Cloning (BC) and Adversarial (ADV) proxies on the performance relative to the expert and the equilibrium imitation learning problem. To this end, we first compute a Nash equilibrium, designated as the expert policy $\pi^E$. 
The objective of the learner is twofold: to obtain a high reward by matching the expert's performance and to learn a Nash equilibrium policy. We compare the case where the learner explicitly accounts for common noise or not.
Thus, we derive two policies from agent trajectories generated by this expert: {\bf (1)}
a population-\textit{un}aware \defi{vanilla policy} $\hat{\pi}^{\text{vanilla}}$, and {\bf (2)} a population-\textit{aware} \defi{adaptive policy} $\hat{\pi}^{\text{adaptive}}$; so $\hat{\pi}^{\text{vanilla}}_n:\cX \to \Delta_{\cA}$ while $\hat{\pi}^{\text{adaptive}}_n: \cX \times \Delta_{\cX} \to \Delta_{\cA}$. We propose two methods to obtain these policies. First, a method well-suited for small spaces and based on Nadaraya-Watson kernel regression. Second, an interactive imitation approach, where a single learning agent is placed within an environment populated by other agents following the expert policy. In all the examples, the proxies will be computed via a Monte Carlo procedure described in Algorithm~\ref{algo:proxy} in Appx.~\ref{appendix:algo}.



\subsection{Challenges}

\noindent\textbf{Nash Equilibrium:}
Computing Nash equilibrium for MFGs without common noise is a well-understood problem in the literature~\cite{achdou2020mean}. However, a gap remains concerning MFGs with common noise: Most existing works consider common noise path-dependent policies, e.g. \cite{achdou2018mean,perrin_master_2022}. However, this becomes intractable as the dimension of the common noise and the time horizon increase.

% In our first example, we utilize an ad-hoc method justified by the low exploitability achieved. For the subsequent two examples, we propose a novel algorithm based on \textbf{Fictitious Play (FP)}  for master policies \cite{perrin_master_2022}, adapted here for the common noise setting. Adapting the FP algorithm is, however, insufficient on its own, as standard FP yields a distribution (or sequence) of policies rather than a unique, deployable policy. To bridge this gap, we propose a new deep learning algorithm designed to learn how to reproduce the trajectory of a mean-field flow. 

% While a theoretical study of these algorithms is beyond the scope of this paper, we highlight their novelty as a potential foundation for future research in the field.

\noindent\textbf{Imitation Learning:} \citep{Chen2022IndividualIRLMFG,Anahtarci2025MaxCausalEntropyMFG,Zhao2025MeanFieldCorrelatedIL} proposed methods to learn policies for MFGs without common noise, but there is no such work with common noise. Here, we need to deal with stochastic mean fields and population-dependent policies. %Here, we highlight a novel learning paradigm: a single agent interacting within an environment populated by other agents already following a Nash equilibrium expert policy. This setting, representing an agent's adaptation to a pre-existing collective behavior under common noise, constitutes a unique research problem that could be studied on its own. We propose a deep imitation learning algorithm specifically for this context.

\subsection{Environment and Metrics}
We consider three different examples with different effects of common noise. The first example is a simple two-state, two-action space. The second example is the Beach Bar environment, and the last example is the Night Club environment inspired by a social choice model~\citep{salhab2019collective}. 

The motivation for testing on three different environments is to demonstrate that our approach is efficient in a variety of structures of common noise:
\begin{itemize}%[noitemsep, topsep=2pt]

    \item \textbf{Exp. 1, Concentration-Dependent:}
    The common noise is correlated with the population density
    (e.g., if the population is concentrated at a state, agents are more likely to be affected).

    \item \textbf{Exp. 2, Local Effects:}
    The noise affects agent positions heterogeneously, leading to more chaotic local shifts.

    \item \textbf{Exp. 3, Block-Shift:}
    The common noise shifts specific clusters or ``blocks'' of agents while leaving others unaffected, representing a segmented environmental shock.

\end{itemize}
In real-world applications, common noise is complex and difficult to model explicitly. By showing consistent results across these three distinct scenarios, we validate that our population-aware imitation learning method remains superior regardless of the specific common noise manifestation.

After computing an expert policy $\pi^E$, a candidate policy $\pi^A$ (which can be either the vanilla policy $\hat \pi^{\text{vanilla}}$ or the adaptive policy $\hat \pi^{\text{adaptive}}$), we report six metrics:
\begin{itemize}%[noitemsep,topsep=0pt]
    \item \textbf{Imitation Proxies:} Behavioral Cloning (BC) and Adversarial (ADV) errors.
    \item \textbf{Performance vs. Expert:} Reward vs. Expert, and Relative Reward vs. Expert: $\cV(\pi^A, \pi^E) = [V(\pi^A, \pi^E) - V(\pi^E, \pi^E)] /\vert V(\pi^E, \pi^E)\vert$.
    \item \textbf{Equilibrium Quality:} Exploitability, and Relative Exploitability: $\cE(\pi^A, \pi^E) = [\cE(\pi^A)]/ \vert V(\pi^E, \pi^E)\vert$.
\end{itemize}
We normalize reward and exploitability metrics to obtain quantities that can be compared across environments. 

To ensure statistical stability, for each parameter configuration, we perform $5$ independent runs of the full experimental pipeline: computing the expert equilibrium, estimating the vanilla and adaptive policies, and calculating the performance metrics. 


\subsection{Fictitious Play for MFGs with Common Noise}
\label{subsection:nash_equilibrium_cn}

As noted earlier, the literature on algorithms for computing Nash equilibria with population-dependent policies for MFGs with common noise is rather scarce. Here we adapt to situations with common noise the Fictitious Play algorithm proposed in~\cite{perrin_master_2022} to obtain a population-dependent Nash equilibrium policy referred to as a \defi{Master policy} (without common noise). 
%In the absence of common noise, Master policies represent the class of population-dependent policies, which are particularly useful for finding equilibria when the initial distribution $\rho_0$ is random rather than fixed. 
% 
% Given that our policies share the same structure as Master policies, we generalize the Fictitious Play (FP) framework to our common noise setting. In FP for MFGs, 
At each iteration, a best response is trained against an aggregated population distribution generated by the sequence of policies obtained in previous iterations. 
%See details in Appx.~\ref{app:FP-algo}.

Specifically, given the policies $(\pi^0, \dots, \pi^k) \in \Pi^{k+1}$ up to iteration $k$, 
% consider the aggregated mean-field sequence $\bar{\rho}^k = (\rho_0, \bar{\rho}^k_1, \dots, \bar{\rho}^k_{H-1})$, where for each time $n \in \{1, \dots, H-1\}$ and common noise history $\varepsilon^0_{1:n}$, the distribution is given by:
% $$
% \bar{\rho}^k_{n}(\varepsilon^0_{1:n}) = \frac{1}{k+1}\sum_{i=0}^k \rho^{\pi^i}_n(\varepsilon^0_{1:n}).
% $$
a policy $\pi^{k+1}$ is defined as a best response to the sequence $(\pi^0, \dots, \pi^k)$ if it maximizes:
% \vspace{-15pt}
$$
    \pi \mapsto \mathbb{E} \Big[ \sum_{n=0}^{H-1} r(X_n, \alpha_n, \bar{\rho}^{k, \varepsilon^0}_n) \Big],
% \vspace{-11pt}
$$
where $(X_n, \alpha_n)$ is induced by $\pi$ and the \defi{aggregate mean-field} sequence $\bar{\rho}^k$ is defined as
% \vspace{-17pt}
\begin{equation}
    \label{eq:aggregate-mf}
    \bar{\rho}^{k, \varepsilon^0}_{n} = \frac{1}{k+1}\sum_{i=0}^k \rho^{\pi^i, \varepsilon^0}_n.
    % \vspace{-12pt}
\end{equation}
The generalized Fictitious Play procedure is summarized in Alg.~\ref{algo:fictitious}.
\begin{algorithm}[t]
\caption{Fictitious Play for MFG with Common Noise}
\begin{algorithmic}[1]
\label{algo:fictitious}
\REQUIRE MFG $\cM = (\cX, \cA, \cE^0, \nu^0, P, r, H, \rho_0)$; Initial policy $\pi^0 \in \Pi$; Number of iterations $K$.
\FOR{$k = 1, \dots, K$}
    \STATE Train $\pi^k$ as the best response to $(\pi^0, \dots, \pi^{k-1})$.
\ENDFOR
\STATE \textbf{Return} policy tuple $(\pi^0, \dots, \pi^K)$.
\end{algorithmic}
\end{algorithm}

To obtain the best response at each iteration, we parametrize the policy by a neural network (NN), and train the NN using Alg.~\ref{algo:nn_best_response} 
% \gl{I use this one only for the ad-hoc method. It is Alg.~\ref{algo:nn_best_response} that i use for beach bar and night club.}
% based on backward induction, 
by directly maximizing the value function, see Appx.~\ref{appendix:algo}. Note that Alg.~\ref{algo:fictitious} returns a list of policies instead of a single, equilibrium policy. To compute $\bar\rho^K$, one can average results obtained with each policy $\pi^k$ of the list, but this is not sufficient for our purpose since we want to have in hand an expert policy $\pi^E$. 
% and the equilibrium it produces is not straightforwardly a Nash equilibrium in the classical sense. Indeed, it is the \textit{aggregated population distribution} that is considered the equilibrium. 
% However, when the transition dynamics $P$ depend on the population distribution, it is not evident that there exists a policy capable of generating such a mean-field sequence. This representation problem is frequently forgotten in the existing literature. 
% For our purposes, this poses a crucial challenge: to evaluate Behavioral Cloning (BC) and Adversarial (ADV) proxies, we require a direct comparison between two individual policies. 
To address this, we propose an IL algorithm that trains a NN to reproduce the aggregated mean-field trajectories, as detailed in Alg.~\ref{algo:il_agg_meanfield}.
\begin{algorithm}[t]
\caption{Mean Field Imitation Learning (MF-IL)}
\begin{algorithmic}[1]
\label{algo:il_agg_meanfield}
\REQUIRE MFG $\cM$; Policy sequence $(\pi^0, \dots, \pi^K)$; 
Initial parameters $\theta_0$ for parametric policy $\pi_\theta$; %Optimizer; 
Iteration number $J$; Batch size $B$.
\FOR{$j = 1, \dots, J-1$}
    % \STATE Reset batch loss: $L(\theta_{j-1}) = 0$.
    \STATE Sample common noises $(\varepsilon^{0, s})_{1\leq s\leq B}$. %a batch of i.i.d. common noise sequences $(\varepsilon^{0, s}_t)_{\substack{1\leq s\leq B \\ 0\leq t \leq H-1}}$.
    \FOR{$s = 1, \dots, B$}
        % \STATE Sample common noises $(\varepsilon^{0, s}_t)_{0\leq t \leq H-1}$.
        \STATE Compute $\bar{\rho}^{K,\varepsilon^{0, s}}$ from $(\pi^0, \dots, \pi^K)$, see~\eqref{eq:aggregate-mf}
        \STATE Generate $\rho^{\pi_{\theta_{j-1}},\varepsilon^{0, s}}$ induced by $\pi_{\theta_{j-1}}$.
    \ENDFOR
    \STATE Compute 
        $L(\theta_{j-1}) = \frac{1}{B} \sum_{s,t} \|\rho^{\pi_{\theta_{j-1}},\varepsilon^{0, s}}_t - \bar{\rho}^{K,\varepsilon^{0, s}}_t \|_1.$
    \STATE Obtain $\theta_j$ by one step of gradient of $L(\theta_{j-1})$.
\ENDFOR
\STATE \textbf{Return} Trained parametric policy $\pi_{\theta_J}$.
\end{algorithmic}
\end{algorithm}

\subsection{Learning Algorithms}
\label{section:learning_algorithm}
To demonstrate the relevance of BC and ADV divergences in distinguishing between population-aware and population-unaware policies, we propose a framework capable of learning both policy types. 
%While the reconstruction of policies from an observed Nash equilibrium in MFG has been explored in existing literature \cite{}, there remains a significant gap regarding games endowed with common noise.

Since our primary objective is to illustrate the efficacy of BC and ADV metrics, we introduce a simple IL algorithm, see Alg.~\ref{algo:interactive_imitation}. This setup mirrors a scenario where a single agent interacts with an environment populated by expert agents who are already following a Nash equilibrium policy $\pi^E$. We assume the learner can observe the states and actions of these expert agents, allowing it to internalize the relationship between the common-noise-induced mean-field $\rho$ and the resulting expert actions.
For population-unaware policies, the same algorithm can be used but removing the mean-field input in $\pi^{\theta_j}$. 

\begin{algorithm}[t]
\caption{Interactive Imitation Learning for Population-aware Policies}
\begin{algorithmic}[1]
\label{algo:interactive_imitation}
\REQUIRE MFG $\cM = (\cX, \cA, \cE^0, \nu^0, P, r, H, \rho_0)$; 
Expert policy $\pi^E$; Initial parameters $\theta_0$; 
Iteration number $J$; Batch size $B$; Population size $M$. \\ 
\FOR{$j = 1, \dots, J$}
    % \STATE Reset batch loss: $L(\theta_{j-1}) = 0$.
    \STATE Sample common noise $(\varepsilon^{0, s}_t)_{0\leq t \leq H-1}^{1\leq s\leq B}$.
    \FOR{$s = 1, \dots, B$}
        \STATE Generate $M$ state-action trajectories $\{(x^{s,i}_t, a^{s,i}_t)\}_{i=1}^M$ using $\pi^E$ under noise $\varepsilon^{0,s}_{1:H-1}$.
        \STATE Set empirical MF: $\hat{\rho}^{s, M}_t(x) = \frac{1}{M}\sum_{i=1}^M \delta_{\{x = x^{s,i}_t\}}$.
        \STATE Set empirical strategy: $\hat{\pi}^{s, M}_t(a \vert x) = \frac{1}{M \hat{\rho}^{s, M}_t(x)} \sum_{i=1}^M \delta_{\{(x, a) = (x^{s,i}_t, a^{s,i}_t)\}}$.
    \ENDFOR
    \STATE Define loss: 
        $
            L(\theta_{j-1}) =
            \frac{1}{BM} 
            % \sum_{s=1}^B\sum_{t=0}^{H-1} \sum_{i=1}^{M} 
            \sum_{s,t,i}\| \pi^{\theta_{j-1}}_t(x^{s,i}_t, \hat{\rho}^{s, M}_t) - \hat{\pi}_t^{s,M}(x^{s,i}_t) \|_1.
        $
    \STATE Obtain $\theta_j$ by one step of gradient of $L(\theta_{j-1})$.
\ENDFOR
\STATE \textbf{Return} Trained adaptive policy $\pi_{\theta_J}$.
\end{algorithmic}
\end{algorithm}
% \vspace{-8pt}
\subsection{First Experiment: 2 States, 2 Actions Environment}
\label{subsection:example1}
\noindent\textbf{Model:} 
Let $\cX = \{0, 1\}$, $\cA = \{0 ,1\}$, $\cE^0 = [0, 1]$, and initial distribution $\rho_0 = (\frac{1}{2}, \frac{1}{2})$. The dynamics are as follows:
For any $e^0\in [0, 1]$ and $\rho \in \Delta_\cX$, we introduce the \emph{perturbed} probability measure $[e^0\rho]$ defined by:
$
    [e^0\rho](1) = \frac{e^0\rho(1)}{e^0\rho(1) + (1 - e^0)\rho(0)}.
$
For any $\eta \in [0, 1]$ and for all $(x,a, \rho, e^0) \in \cX \times \cA \times \Delta_\cX\times [0, 1]$,
$
        P_\eta(\cdot\vert x, a,\rho, e^0) = (1 - \eta)\delta_a + \eta [e^0\rho].
$
The reward is:
$
r(x,a, \rho) = -\mathds{1}_{\{x = 0\}}\rho(0)  -\mathds{1}_{\{x = 1\}}\rho(1).
$
For the common noise distribution, 
we consider $\nu^0_{\alpha} = \texttt{Beta}(\alpha, \alpha)$. 
%
This model represents agents who choose between two alternatives under congestion effects. Individual decisions are occasionally ignored and replaced by a population-level force. When this force is active, its impact is governed by $\alpha$. When $\alpha < 1$, $\texttt{Beta}(\alpha, \alpha)$ is U-shaped and $ [\varepsilon^0_1\rho]$ concentrates at $0$ or $1$, with equal probability. When $\alpha =1$, $\texttt{Beta}(1, 1)$ is the uniform law, and the behavior of  $ [\varepsilon_1^0\rho]$ is erratic.  When $\alpha > 1$, $\texttt{Beta}(\alpha, \alpha)$ is concentrated around $\frac{1}{2}$ and the perturbation $ [\varepsilon^0_1\rho]$ stays close to $\rho$. In short, $\eta$ controls the frequency at which the common noise impacts the agents, and $\alpha$ controls its effect.

\noindent\textbf{Expert Policy:} 
To compute a Nash equilibrium on this simple environment, we use an ad-hoc method inspired by Mann iterations 
% \ml{not sure if everyone is familiar, and Mann iteration for which function $f$? BR? But it is multi-valued} 
% \gl{For BR with DPP. But we can remove the term 'Mann' and just say 'ad-hoc'. } 
of best responses. The best response is computed via dynamic programming following \citep{gu2024global}, adapted to incorporate common noise.  
This method is made possible by the simplicity of the environment; details of the algorithm are presented in Appx.~\ref{appendix:example1}. The relevance of this method is justified by the low exploitability achieved by the expert policy. For each tuple of parameters $(\alpha, \eta) \in \mathbb{R}_+ \times [0, 1]$, we denote the expert policy as $\pi^E_{\alpha, \eta}$.

\noindent\textbf{Vanilla and Adaptive Policies:}
The simplicity of the environment allows us to consider classical statistical estimators for the vanilla and adaptive policies. For this example, we propose a method based on Nadaraya-Watson kernel regression. To this end, we sample synthetic state-action trajectories of agents following the expert policy under the same common noise realizations. This allows us to regress the expert policy by observing the actions alongside the empirical population distribution. The adaptive policy thus interpolates the expert's behavior across the continuous space of mean-field realizations, while the vanilla policy serves as a baseline that ignores common noise fluctuations. Complete technical details are provided in Appx.~\ref{appendix:example1}.

\noindent\textbf{Results:} We consider $\alpha \in \{0.75, 1, 1.25, 1.5, 1.75\}$ and $\eta \in \{0, 0.25, 0.5, 0.75, 1\}$. 
% To ensure statistical stability, for each parameter configuration $(\alpha, \eta)$, we perform $5$ independent runs of the full experimental pipeline: computing the expert equilibrium, estimating the vanilla and adaptive policies, and calculating the performance metrics. 

The complete grid of results is provided in Appx.~\ref{appendix:example1}. We highlight specific findings for $\eta = 0.75$ in Figure~\ref{fig:example1_eta=0.75} and for $\alpha = 1.75$ in Figure~\ref{fig:example1_sub_alpha} in Appx.~\ref{appendix:example1}. The results indicate that the problem exhibits significant parameter sensitivity. Notably, the adaptive policy consistently outperforms the vanilla policy across all tested regimes. 

This performance gap can be explained by a qualitative analysis of the learned policies. Figure~\ref{fig:example1_behavior} illustrates the behavior of $\pi^E, \hat{\pi}^{\text{vanilla}}$, and $\hat{\pi}^{\text{adaptive}}$ as a function of the mean-field variable $\rho$. While the expert dynamically adapts its actions to the fluctuations of the mean-field, the adaptive policy successfully reconstructs this dependency. In contrast, the vanilla policy only captures the average action of the expert.
% \begin{figure}[htbp]
%     \centering
%     \includegraphics[width=0.9\linewidth]{Figures/Example1/ex1_eta=0.75_QUAD.pdf}
%     \caption{Performance metrics for $5$ different runs, with $\eta = 0.75$.}
%     \label{fig:example1_eta=0.75}
% \end{figure}
% \begin{figure}[htbp]
%     \centering
%     \includegraphics[width=0.9\linewidth]{Figures/Example1/Ex1_t=10_alpha=1.75_eta=0.75_behavior.pdf}
%     \caption{Policies comparison for $\alpha = 1.75$, $\eta = 0.75$.}
%     \label{fig:example1_behavior}
% \end{figure}


\begin{figure}[htbp] % Recommended for top of page
    \centering
    \includegraphics[width=0.9\linewidth]{Figures/Example1/ex1_eta=0.75_QUAD.pdf}
    \caption{Performance metrics for $5$ different runs, with $\eta = 0.75$.}
    \label{fig:example1_eta=0.75}

    \vspace{0.3cm} % Explicitly control the gap here

    \includegraphics[width=0.9\linewidth]{Figures/Example1/Ex1_t=10_alpha=1.75_eta=0.75_behavior.pdf}
    \caption{Policy comparison for $\alpha = 1.75$, $\eta = 0.75$.}
    \label{fig:example1_behavior}
    % \vspace{-5pt}
\end{figure}

\subsection{Second Experiment: The Beach Bar Environment }
We consider the \textit{Beach Bar} environment as presented in \cite{perrin_master_2022}. In this setting, the bar is permanently open and the agents are subject to a common noise. 

\noindent\textbf{Model:}
Let $\cX = \{0, \dots, 4X\}$, $\cA = \{-1, 0, 1\}$, and $\cE^0 = \cX^{|\cX|}$.
For parameters $\alpha, \beta, \eta$, and common noise $e^0 \in \cE^0$, the dynamics are governed by a transition kernel $P$ that incorporates both \emph{common noise} $e^0(x)$ and \emph{idiosyncratic noise} (represented by a uniform shift $s \in \{-1, 0, 1\}$). Specifically,
$P(x, a, \rho, e^0) = \frac{1}{3} \sum_{s \in \{-1, 0, 1\}} \delta_{x + a + e^0(x) + s \pmod{|\cX|}}$.
The common noise law is defined as $\varepsilon^0_1 \sim (\nu^0_\eta)^{\otimes |\cX|}$, where $\nu^0_\eta$ is the distribution of $\varepsilon^{0,1} \varepsilon^{0,2}$ with $\varepsilon^{0, 1} \sim \text{Bernoulli}(\eta)$ and $\varepsilon^{0, 2} \sim \text{Unif}\{-X, \dots, X\}$.
The reward function, reflecting a preference for a central bar at $x_{\text{bar}} = 2X$ and a penalty for overcrowding, is:
$r(x, a, \rho) = -|x - x_{\text{bar}}| \mathds{1}_{\{\rho(x_\text{bar}) \leq \beta\}} - \alpha \log(\rho(x)) - |a|$.
This model represents a beach environment where holidaymakers balance two competing objectives: avoiding crowds while remaining close to the bar. However, the agents are subject to a threshold effect; when the bar becomes too crowded ($\rho(x_{\text{bar}}) > \beta$), it loses its appeal, and agents lose all interest to it. The environment is disrupted by external events that simultaneously displace all agents at specific locations. These may represent physical changes, such as shifting shadows making certain spots undesirable, or social phenomena like spontaneous group movements. These stochastic events are governed by the common noise, where the parameter $\eta$ controls the frequency of such perturbations.

For the following, we fix $\beta = 0.1$. We study the effect of the parameters $(\alpha, \eta)$ on our learning problems. 

\noindent\textbf{Expert Policy:}
For this large scale environment, we use our new algorithmic procedure described in Subsec.~\ref{subsection:nash_equilibrium_cn}. Figure~\ref{fig:beachbar_meanfield} shows one realization of the mean-field trajectory induced by the expert, for $(\alpha, \eta) = (1, 0.3)$.
\begin{figure}[htbp]
    \centering
    \includegraphics[width=0.9\linewidth]{Figures/BeachBar/beach_bar_eta=0.3_alpha_1_meanfield.pdf}
    \caption{One realization of the mean-field trajectory generated by the expert, for $(\alpha, \eta) = (1, 0.3)$.}
    \label{fig:beachbar_meanfield}
    % \vspace{-5pt}
\end{figure}

\noindent\textbf{Vanilla and Adaptive Policies:}
For each value of $(\alpha, \eta)$, we consider learning algorithms presented in Subsec.~\ref{section:learning_algorithm}. We report all the learning parameters used in Appx.~\ref{appendix:beachbar}.

\noindent\textbf{Convergence of the Algorithms:}
Figure~\ref{fig:beachbar_losses_eta=0.3_alpha=1} reports, for $(\alpha, \eta) = (1, 0.3)$, the convergence of the FP procedure, the expert loss from Algorithm~\ref{algo:il_agg_meanfield}, and the losses of the adaptive and vanilla policies from Algorithm~\ref{algo:interactive_imitation}. The FP procedure is confirmed to converge toward a Nash equilibrium as the number of iterations grows. Moreover, the adaptive policy achieves lower loss than the vanilla policy, indicating superior imitation of the expert. Additional results for $\eta = 0.3$ across multiple values of $\alpha$ are provided in Appx.~\ref{appendix:beachbar}, Figure~\ref{fig:beachbar_losses}.


\vspace{1pt}
\begin{figure}[htbp]
    \centering
    \includegraphics[width=1\linewidth]{Figures/BeachBar/beachbar_losses_eta=0.3_alpha=1_eta0.3_alpha1.pdf}
    \caption{Convergence curves for $(\alpha, \eta) = (1, 0.3)$: FP exploitability (left), expert loss (center), and adaptive vs.\ vanilla losses (right).}
    \label{fig:beachbar_losses_eta=0.3_alpha=1}
    % \vspace{-5pt}
\end{figure}
\noindent\textbf{Results:}
We take $\eta \in \{0.1, 0.2, 0.3, 0.4, 0.5\}$ and $\alpha \in \{0.1,0.5,1,1.5,2\}$, $X = 5$ ($\vert \cX\vert = 20)$ and $H = 50$. 
% To ensure statistical stability, for each parameter configuration $(\alpha, \eta)$, we perform $5$ independent runs of the full experimental pipeline: computing the expert equilibrium, estimating the vanilla and adaptive policies, and calculating the performance metrics.

The complete grid of results is provided in Appx.~\ref{appendix:beachbar}. We highlight specific findings for  $\eta = 0.3$ in Figure~\ref{fig:beachbar_eta=0.3} and $\alpha = 1$ in Figure~\ref{fig:beachbar_sub_alpha} in Appx.~\ref{appendix:beachbar}. The results indicate that the problem exhibits significant parameter sensitivity. We again observe that the adaptive policy consistently outperforms the vanilla policy across all tested regimes. 

The qualitative analysis for this environment mirrors that of Example 1. The adaptive policy successfully reconstructs the expert’s dependency on $\rho$, whereas the vanilla policy fails to react to mean-field fluctuations, as we can see in Figure~\ref{fig:beachbar_policies}. This explains the observed performance disparity.

% \begin{figure}[htbp]
%     \centering
%     \includegraphics[width=0.9\linewidth]{Figures/BeachBar/beachbar_eta=0.3_FULL.pdf}
%     \caption{ Performance metrics for 5 different runs, with $\eta = 0.3$.}
%     \label{fig:beachbar_eta=0.3}
% \end{figure}
% \begin{figure}[htbp]
%     \centering
%     \includegraphics[width=0.9\linewidth]{Figures/BeachBar/behavior_meanfield.pdf}
%     \caption{Policy sensitivity analysis for the Beach Bar environment with parameters $(\alpha, \eta) = (1, 0.3)$ at time $n = 10$ and state $x = x_{\text{bar}}$. The comparison illustrates how the Expert (Nash), Adaptive, and Vanilla policies respond to different population distributions:\textit{Crowded Bar} represents high density at the bar location, testing the agent's avoidance behavior; \textit{Shifted Down} and \textit{Shifted Up} reflect scenarios where the mass is concentrated immediately below or above the bar, respectively. This demonstrates the Adaptive policy's ability to react to mean-field fluctuations, whereas the Vanilla policy remains static due to its lack of population awareness.}
%     \label{fig:beachbar_policies}
% \end{figure}
\begin{figure}[!t] % Recommended for top of page
    % \vspace{-8pt}
    \centering
    \includegraphics[width=0.99\linewidth]{Figures/BeachBar/beachbar_eta=0.3_QUAD.pdf}
    \caption{Performance metrics for 5 different runs, with $\eta = 0.3$.}
    \label{fig:beachbar_eta=0.3}

    \vspace{0.3cm} % Explicitly control the gap here

    \includegraphics[width=0.99\linewidth]{Figures/BeachBar/behavior_meanfield.pdf}
    \caption{Policy sensitivity analysis for the Beach Bar environment with parameters $(\alpha, \eta) = (1, 0.3)$ at time $n = 10$ and state $x = x_{\text{bar}}$. The comparison illustrates how the Expert (Nash), Adaptive, and Vanilla policies respond to different population distributions:\textit{Crowded Bar} represents high density at the bar location, testing the agent's avoidance behavior; \textit{Shifted Down} and \textit{Shifted Up} reflect scenarios where the mass is concentrated immediately below or above the bar, respectively. This demonstrates the Adaptive policy's ability to react to mean-field fluctuations, whereas the Vanilla policy remains static due to its lack of population awareness.}
    \label{fig:beachbar_policies}
    % \vspace{-10pt}
\end{figure}
% \vspace{-10pt}
\subsection{Third Experiment: The Night Clubs Environment}
The following model is inspired by the mean-field linear-quadratic-Gaussian (LQG) model for social choice developed by~\citet{salhab2019collective}. This environment shares structural similarities with the \textit{Beach Bar} problem. However, the reward structure is different: the environment features two distinct points of interest (the "night clubs"), and agents are motivated by a gathering incentive, a desire to congregate rather than avoid congestion. Furthermore, the common noise in this setting is designed to affect agents in a more highly correlated manner than in previous experiments, simulating global trends that shift the popularity of one club over another.

\noindent\textbf{Model:}
Let $\cX, \cA, \cE^0$, and dynamics $P$ be as defined in the Beach Bar model, with clubs at $x^1_{\text{club}}=X$ and $x^2_{\text{club}}=3X$. The reward function $r(x, a, \rho) = -D_{\beta}^k(x, \rho) - \alpha |x - \bar{x}(\rho)|$ prioritizes a target club based on crowding and penalized distance from the circular mean $\bar{x}(\rho)$. The effective distance is $D_{\beta}^k(x, \rho) = |x - x^k_{\text{club}}|$ if $\rho(x^k_{\text{club}}) \le \beta$, and $|x - x^{-k}_{\text{club}}|$ otherwise.
Common noise $\varepsilon^0_1$ follows a localized \textbf{block-shift} law: at each step, a shift magnitude $s \sim \text{Unif}\{-\lfloor \frac{|\cX|}{3} \rfloor, \dots, \lfloor \frac{|\cX|}{3} \rfloor\}$ and a block start $i \sim \text{Unif}(\cX)$ are sampled. The resulting noise vector is $\varepsilon^0_1(x) = s \cdot \mathds{1}_{\{x \in \{i, i+1 \pmod{|\cX|}\}\}} \cdot \eta^{1/3}$, effectively shifting only a local block of agents. This model represents a social environment where clubgoers balance two competing objectives: the desire for social gathering and the pursuit of entertainment at local venues. Agents seek to minimize their distance to the nearest club to enjoy the music; however, the clubs are subject to a mechanical threshold effect. If a club's local density exceeds the capacity $\beta$ ($\rho(x_{\text{club}}) > \beta$), the venue is considered overcrowded, the music ceases, and the agents lose all incentive to remain at that location. Agents immediately redirect their interest toward the alternative club. This redirection is 'blind', the agents are assumed to be too distant to perceive the current status (open or closed) of the far venue, and thus target its location regardless of its density. The environment is disrupted by structured external events that simultaneously displace groups of agents, modeling sudden interest in the direction of one of the music venues.
For the following, we fix $\beta = 0.1$. We study the effect of the parameters $(\alpha, \eta)$ on our learning problems. 

\noindent\textbf{Expert, Vanilla and Adaptive Policies:}
We consider the exact same procedure as in the Beach Bar example. Hyperparameter details and training loss are given in Appx.~\ref{appendix:nightclubs}. Figure~\ref{fig:nightclub_meanfield} shows one realization of the mean-field trajectory induced by the expert, for $(\alpha, \eta) = (0.1, 1)$. Figure~\ref{fig:nightclubs_losses_eta=1_alpha=0.1} reports the convergence of the FP procedure and the imitation losses for the same parameter configuration. As in the Beach Bar experiment, the FP procedure converges toward a Nash equilibrium, and the adaptive policy achieves lower loss than the vanilla policy, confirming its superior imitation of the expert.

\begin{figure}[htbp]
    \centering
    \includegraphics[width=0.9\linewidth]{Figures/NightClubs/nightclub_meanfield_meanfield.pdf}
    \caption{One realization of the mean-field trajectory generated by the expert for $(\alpha, \eta) = (0.1, 1)$.}
    \label{fig:nightclub_meanfield}
\end{figure}


% \vspace{5pt}
\begin{figure}[htbp]
    \centering
    \includegraphics[width=1\linewidth]{Figures/NightClubs/nightclub_losses_eta=1_alpha=0.1_eta1_alpha0.1.pdf}
    \caption{Convergence curves for $(\alpha, \eta) = (0.1, 1)$: FP exploitability (left), expert loss (center), and adaptive vs.\ vanilla losses (right).}
    \label{fig:nightclubs_losses_eta=1_alpha=0.1}
\end{figure}

% \begin{figure}[htbp]%[H] % Recommended for top of page
%     \centering
%     \includegraphics[width=0.9\linewidth]{Figures/NightClubs/nightclub__alpha=1.05_QUAD.pdf}
%     \caption{ Performance metrics for 5 different runs, with $\alpha = 1.05$.}
%     \label{fig:nightclubs_alpha=1.05}

%     % \vspace{0.3cm} % Explicitly control the gap here

%     \includegraphics[width=0.85\linewidth]{Figures/NightClubs/behavior_meanfield.pdf}
%     \caption{Policy sensitivity analysis for the Night Clubs environment with parameters $(\alpha, \eta) = (0.1, 1)$ at time $n = 10$ and state $x = 2X$. The comparison illustrates how the Expert (Nash), Adaptive, and Vanilla policies respond to different population distributions:\textit{Crowded Club 1} represents high density at the Club 1 location; \textit{Balanced Clubs} represents equal density at Club 1 and Club 2 locations. \textit{Crowded Club 2} represents high density at the Club 2 location.}
%     \label{fig:nightclubs_policies}
% \end{figure}

\begin{figure}[H]
    \centering
    \includegraphics[width=0.99\linewidth]{Figures/NightClubs/nightclub__alpha=1.05_QUAD.pdf}
    \caption{Performance metrics for 5 different runs, with $\alpha = 1.05$.}
    \label{fig:nightclubs_alpha=1.05}
\end{figure}

\begin{figure}[H]
    \centering
    \includegraphics[width=0.99\linewidth]{Figures/NightClubs/behavior_meanfield.pdf}
    \caption{Policy sensitivity analysis for the Night Clubs environment with parameters $(\alpha, \eta) = (0.1, 1)$ at time $n = 10$ and state $x = 2X$. The comparison illustrates how the Expert (Nash), Adaptive, and Vanilla policies respond to different population distributions: \textit{Crowded Club 1} represents high density at the Club 1 location; \textit{Balanced Clubs} represents equal density at Club 1 and Club 2 locations. \textit{Crowded Club 2} represents high density at the Club 2 location.}
    \label{fig:nightclubs_policies}
\end{figure}
\noindent\textbf{Results:}
We take $\eta \in \{0.1, 0.325, 0.55, 0.775, 1\}$ and $\alpha \in \{0.1,0.575,1.05,1.525,2\}$, $X = 5$ ($\vert \cX\vert = 20)$ and $H = 30$.
% To ensure statistical stability, for each parameter configuration $(\alpha, \eta)$, we perform $5$ independent runs of the full experimental pipeline.




The complete grid of results is provided in Appx.~\ref{appendix:nightclubs}. We highlight specific findings for $\alpha = 1.05$ in Figure~\ref{fig:nightclubs_alpha=1.05} and for $\alpha = 0.1$ in Figure~\ref{fig:nightclubs_sub_alpha=0.1} in Appx.~\ref{appendix:nightclubs}. The results indicate that the problem exhibits significant parameter sensitivity. We again observe that the adaptive policy consistently outperforms the vanilla policy across all tested regimes, but that the advantage on \ref{theory-q2} (reward matching) is less obvious. 

The qualitative behavior in the Night Clubs environment mirrors our previous observations. As shown in Figure~\ref{fig:nightclubs_policies}, the adaptive policy effectively reconstructs the expert’s reactivity to mean-field fluctuations $\rho$, whereas the vanilla policy fails to capture these dynamics, serving as a mere average of expert actions.
% \begin{figure}[htbp]
%     \centering
%     \includegraphics[width=0.9\linewidth]{Figures/NightClubs/nightclub__alpha=1.05_FULL.pdf}
%     \caption{ Performance metrics for 5 different runs, with $\alpha = 1.05$.}
%     \label{fig:nightclubs_alpha=1.05}
% \end{figure}
% \begin{figure}[htbp]
%     \centering
%     \includegraphics[width=0.9\linewidth]{Figures/NightClubs/behavior_meanfield.pdf}
%     \caption{Policy sensitivity analysis for the Night Clubs environment with parameters $(\alpha, \eta) = (0.1, 1)$ at time $n = 10$ and state $x = 2X$. The comparison illustrates how the Expert (Nash), Adaptive, and Vanilla policies respond to different population distributions:\textit{Crowded Club 1} represents high density at the Club 1 location; \textit{Balance Clubs} represents equal density at Club 1 and Club 2 locations.\textit{Crowded Club 2} represents high density at the Club 2 location.}
%     \label{fig:beachbar_policies}
% \end{figure}

% \vspace{-10pt}
\section{Conclusion}
\label{sec:ccl}

We present the first framework for Imitation Learning in Mean Field Games with common noise. By extending Behavioral Cloning and Adversarial proxies to stochastic mean fields, we prove these metrics control both exploitability and performance relative to an expert. Our analysis highlights the critical role of \defi{adaptive policies}, showing that population-unaware strategies fail to capture equilibrium dynamics under stochastic fluctuations. To validate this, we developed a novel pipeline combining generalized Fictitious Play and Deep Learning to approximate population-aware policies. Future work should focus on extending the framework for instance to state-action mean field interactions and applying the method to real-world data.



\section*{Impact Statement}
While our work is primarily theoretical, we believe it could find implications in real-world applications involving massive agent interactions. Due to high dimensionality and the difficulty of characterization, scientific models often omit common noise: the external stochastic signals that affect an entire population simultaneously. However, real-world equilibria are influenced by these global signals. In financial markets, for example, a single statement from a political figure on social media can trigger instantaneous, population-wide shifts in trading volumes and asset prices. Or in energy infrastructure, environmental randomness in renewable energy production affects all producers in a geographic region simultaneously. Our results suggest that in those situations, Imitation Learning with population-aware policies is more robust than population-unaware policies.


% % This paper is the first to address Imitation Learning for Mean Field Games subject to common noise, specifically focusing on the challenges of reconstructing an equilibrium and achieving high rewards against an expert population. We established a comprehensive framework extending the analysis of Behavioral Cloning and Adversarial proxies to stochastic mean field flows. We provided four theoretical results proving that these proxies effectively control the exploitability of the learned policy and its performance gap relative to the expert. A key insight of our work is the necessity of \defi{adaptive policies}: we demonstrated that population-independent ("vanilla") strategies fail to capture equilibrium dynamics when the population distribution fluctuates stochastically. Our numerical results confirm that adaptive policies outperform vanilla policies by accounting for social reactivity.

% Developing our experiments required us to introduce a novel algorithmic pipeline for computing Nash equilibria and designing imitation algorithms in this setting. As these remain open questions, our approach provides a foundation for further research, specifically through our two-stage framework: computing a robust set of policies via generalized Fictitious Play, and subsequently obtaining a policy that reconstructs the aggregated mean field via Deep Learning.

% Future research should continue the theoretical study of this numerical procedure and the proposed Interactive Imitation Learning algorithm. Furthermore, our current model assumes interactions depend solely on the state distribution; extending this to Mean Field Games of controls, where the mean field of \textit{actions} also influences dynamics, would broaden the method's applicability. Finally, the MFG literature would benefit greatly from further illustrations using real-world data to bridge the gap between theory and practice.



% We introduced a comprehensive framework for Imitation Learning in Mean Field Games subject to common noise. We extended the theoretical analysis of Behavioral Cloning and Adversarial proxies to stochastic mean field flows, establishing finite-sample bounds that control both the exploitability of the learned policy and its performance gap relative to the expert. A key insight of our work is the necessity of adaptive policies: we demonstrated that population-independent strategies fail to capture equilibrium dynamics when the population distribution fluctuates stochastically. To support these findings, we developed a novel algorithmic pipeline combining generalized Fictitious Play with Deep Learning to compute and imitate Master policies. 
% %
% Our results open several avenues for future research. First, applying this framework to real-world datasets remains a crucial next step. 
% Second, our current model assumes interactions depend solely on the state distribution. Extending this framework to  Mean Field Games of controls, where the mean field of \textit{actions} also influences dynamics and rewards, would broaden the applicability of our method.


% \noindent\textbf{Impact Statement:} This paper presents work whose goal is to advance the field of Machine Learning. There are many potential societal consequences of our work, none which we feel must be specifically highlighted here.

%  \gl{On ICML web page:
% Authors are required to include a statement of the potential [...] This statement should be in a separate section at the end of the paper (co-located with Acknowledgements, before References), and does not count toward the paper page limit. [...]}
% \gl{It does not count toward the paper page limit, so we should probably say it little more ?? I put my proposition below my conclusino part. }

% \section{Conclusion}

% This paper is the first to address Imitation Learning for MFGs with common noise, specifically focusing on the challenges of reconstructing an equilibrium and achieving high rewards against a expert population. We established four theoretical results proving that Behavioral Cloning (BC) and Adversarial (ADV) proxies effectively control these objectives. To illustrate our results, we analyzed the performance of an agent learning a population-independent (vanilla) policy by observing an expert population, vs the one of an agent learning a population dependent (adaptive) policy. Our numerical results suggest that in the context of multi-agent imitation learning, adaptive policies outperform vanilla policies by accounting for social reactivity.

% Developing our experiments required us to introduce the challenges of computing a Nash equilibrium and designing imitation algorithms specifically for MFGs with common noise. As these remain open questions in the field, our approach provides a foundation for further original research.

% Future research should continue the theoretical study of our numerical procedure to obtain a Nash equilibirum. We emphasize that our current framework consists of two distinct stages: computing a robust set of policies (Fictitious Play algorithm) and subsequently obtaining a policy that reconstructs the aggregated mean field (Mean Field Imitation Algorithm). Additionally, our Interactive Imitation Learning algorithm warrants further investigation. Finally, we note that the MFG literature would benefit greatly from further illustrations using real-world data to bridge the gap between theory and practice.


% While our work is primarily theoretical, its implications are significant for real-world systems involving massive agent interactions, such as crowd dynamics, high-frequency trading, epidemiology, and statistical physics. We emphasize that in multi-agent imitation learning, it is critical to account for population-aware policies. Due to its high dimensionality and difficulty to characterize, scientific models often omit common noise, the external stochastic signals that affect an entire population simultaneously. However, real-world equilibrium are deeply influenced by these global signals.

% To illustrate the complexity and impact of common noise in multi-agent systems, one can consider several domains:
% \begin{itemize}
%     \item Financial Markets: A single statement from a political figure on social media can trigger instantaneous, population-wide shifts in trading volumes and asset prices.
%     \item Epidemiology: Public awareness campaigns during a health crisis act as a common signal that fundamentally alters the contact rates and behavioral patterns of the entire population at once.
%     \item Energy Infrastructure: Environmental alea (randomness) in renewable energy production affects all producers in a geographic region simultaneously, requiring a coordinated equilibrium response to maintain grid stability. 
% \end{itemize}
% By providing theoretical guarantees for policies that adapt to the mean-field, we provide a framework for learning behaviors that remain stable under such noise. Our findings suggest that in real-world applications, practitioners should prioritize learning the social reactivity of a population rather than just individual state-action pairs. Failing to do so could lead to the deployment of "equilibrium-blind" agents that might achieve high local rewards but inadvertently destabilize the global social or economic system when faced with common noise.

\bibliography{biblio}
\bibliographystyle{icml2026}

%%%%%%%%%%%%%%%%%%%%%%%%%%%%%%%%%%%%%%%%%%%%%%%%%%%%%%%%%%%%%%%%%%%%%%%%%%%%%%%
%%%%%%%%%%%%%%%%%%%%%%%%%%%%%%%%%%%%%%%%%%%%%%%%%%%%%%%%%%%%%%%%%%%%%%%%%%%%%%%
% APPENDIX
%%%%%%%%%%%%%%%%%%%%%%%%%%%%%%%%%%%%%%%%%%%%%%%%%%%%%%%%%%%%%%%%%%%%%%%%%%%%%%%
%%%%%%%%%%%%%%%%%%%%%%%%%%%%%%%%%%%%%%%%%%%%%%%%%%%%%%%%%%%%%%%%%%%%%%%%%%%%%%%


\newpage
\appendix
\onecolumn
\allowdisplaybreaks

\section{Useful Inequalities and Proofs}
\label{appendix:theory}

In this appendix, we first present the inequalities that are used to establish our theoretical results. We then proceed to the proofs of the theorems. 

\subsection{Useful Inequalities}

First we extend \citep[Lemma 2 C.4]{ramponi2023imitation}.

\begin{lemma}
    \label{lemma:Vpi3pi1-Vpi3pi2}
    Under Assumptions~\ref{assumption:Lipschitz} and~\ref{assumption:rmax}, for any policies $\pi^1, \pi^2, \pi^3 \in \Pi$, we have:
    \begin{flalign*}
        |V(\pi^3, \pi^1) - V(\pi^3, \pi^2)| &\le L_r \mathbb{E} \left[ \sum_{n=0}^{H-1} \|\rho^1_n - \rho^2_n\|_1 \right] + r_{\max} \mathbb{E} \left[ \sum_{n=0}^{H-1} \|\rho^{1,3}_n - \rho^{2,3}_n\|_1 \right] &\\
        &\quad + r_{\max} \sum_{n=0}^{H-1} \mathbb{E} \left[ \mathbb{E}_{X \sim \rho^{2,3}_n} \left[ \|\pi^3_n(X, \rho^1_n) - \pi^3_n(X, \rho^2_n)\|_1 \right] \right]. &
    \end{flalign*}
\end{lemma}

\begin{proof}
    Expanding the value function difference:
    \begin{flalign*}
        |V(\pi^3, \pi^1) - V(\pi^3, \pi^2)| &= \left| \mathbb{E} \left[ \sum_{n=0}^{H-1} \sum_{x,a} \left( \mu^{1,3}_n(x, a) r(x, a, \rho^{1}_n) - \mu^{2,3}_n(x, a) r(x, a, \rho^{2}_n) \right) \right] \right| &\\
        &\le \underbrace{\mathbb{E} \left[ \sum_{n,x,a} \mu^{1,3}_n(x, a) |r(x, a, \rho^1_n) - r(x, a, \rho^2_n)| \right]}_{A} &\\
        &\quad + \underbrace{\mathbb{E} \left[ \sum_{n,x,a} |(\mu^{1,3}_n(x, a) - \mu^{2,3}_n(x, a)) r(x, a, \rho^2_n)| \right]}_{B}. &
    \end{flalign*}
    
    For term $A$, by the Lipschitz continuity of $r$ in the mean-field:
    \begin{flalign*}
        A &\le \mathbb{E} \left[ \sum_{n=0}^{H-1} L_r \|\rho^1_n - \rho^2_n\|_1 \sum_{x,a} \mu^{1,3}_n(x, a) \right] = L_r \mathbb{E} \left[ \sum_{n=0}^{H-1} \|\rho^1_n - \rho^2_n\|_1 \right]. &
    \end{flalign*}
    
    For term $B$, we use the definition $\mu_n(x,a) = \rho_n(x)\pi_n(a|x, \rho)$ and the triangle inequality:
    \begin{flalign*}
        B &\le r_{\max} \mathbb{E} \left[ \sum_{n,x,a} |\rho^{1,3}_n(x)\pi^3_n(a|x, \rho^1_n) - \rho^{2,3}_n(x)\pi^3_n(a|x, \rho^2_n)| \right] &\\
        &\le r_{\max} \mathbb{E} \left[ \sum_{n,x,a} \pi^3_n(a|x, \rho^1_n) |\rho^{1,3}_n(x) - \rho^{2,3}_n(x)| + \rho^{2,3}_n(x) |\pi^3_n(a|x, \rho^1_n) - \pi^3_n(a|x, \rho^2_n)| \right] &\\
        &= r_{\max} \mathbb{E} \left[ \sum_{n=0}^{H-1} \|\rho^{1,3}_n - \rho^{2,3}_n\|_1 \right] + r_{\max} \sum_{n=0}^{H-1} \mathbb{E} \left[ \mathbb{E}_{X \sim \rho^{2,3}_n} \left[ \|\pi^3_n(X, \rho^1_n) - \pi^3_n(X, \rho^2_n)\|_1 \right] \right]. &
    \end{flalign*}
    This concludes the proof.
\end{proof}

Now we verify \citep[Lemma 3 C.4]{ramponi2023imitation}.
\begin{lemma}
    \label{lemma:rhoE-rhoEA_BC-muE-muEA_BC}
    Recall that $\delta^{BC} = \max_{n \in [H-1]} \EE(\EE_{X \sim \rho^E_n}( \| \pi^E_n(X, \rho^E_n) - \pi^A_n(X, \rho^E_n) \|_1)).$ We have 
    \begin{equation}
        \label{eq:rhoE-rhoEA_BC}
        \EE(\sum_{n=0}^{H-1} \| \rho^{E}_n - \rho^{EA}_n \|_1) \leq H^2 \delta^{BC},
    \end{equation}
    and
    \begin{equation}
        \label{eq:muE-muEA_BC}
        \EE(\sum_{n=0}^{H-1} \| \mu^{E}_n - \mu^{EA}_n \|_1) \leq H^2 \delta^{BC}.
    \end{equation}
\end{lemma}
\begin{proof}
    \begin{align*}
        \EE\!\left(\| \rho^{E}_n - \rho^{EA}_n \|_1\right)
        &= 
        \EE\!\left( \sum_{x'} \lvert \rho^{E}_n(x') - \rho^{EA}_n(x') \rvert \right) \\[4pt]
        &=
        \EE\!\left(
        \sum_{x'} 
        \left\lvert 
        \sum_{x,a} 
        \rho^{E}_{n-1}(x)\,\pi^E_{n-1}(a \vert x,\rho^{E}_{n-1})\,P(x'\vert x,a,\rho^{E}_{n-1}, \varepsilon^0_{n+1})
        -
        \rho^{EA}_{n-1}(x)\,\pi^A_{n-1}(a \vert x,\rho^{E}_{n-1})\,P(x'\vert x,a,\rho^{E}_{n-1}, \varepsilon^0_{n+1})
        \right\rvert
        \right) \\[6pt]
        &=
        \EE\!\left(
        \sum_{x'} 
        \left\lvert 
        \sum_{x,a}
        P(x' \vert x,a,\rho^{E}_{n-1}, \varepsilon^0_{n+1})
        \bigl(
        \rho^{E}_{n-1}(x)\pi^E_{n-1}(a \vert x,\rho^{E}_{n-1})
        -
        \rho^{EA}_{n-1}(x)\pi^A_{n-1}(a \vert x,\rho^{E}_{n-1})
        \bigr)
        \right\rvert
        \right) \\[6pt]
        &\le 
        \EE\!\left(
        \sum_{x'}\sum_{x,a}
        P(x' \vert x,a,\rho^{E}_{n-1}, \varepsilon^0_{n+1})\,
        \left\lvert 
        \rho^{E}_{n-1}(x)\pi^E_{n-1}(a \vert x,\rho^{E}_{n-1})
        -
        \rho^{EA}_{n-1}(x)\pi^A_{n-1}(a \vert x,\rho^{E}_{n-1})
        \right\rvert
        \right) \\[6pt]
        &=
        \EE\!\left(
        \sum_{x,a}
        \left\lvert 
        \rho^{E}_{n-1}(x)\pi^E_{n-1}(a \vert x,\rho^{E}_{n-1})
        -
        \rho^{EA}_{n-1}(x)\pi^A_{n-1}(a \vert x,\rho^{E}_{n-1})
        \right\rvert
        \underbrace{\sum_{x'} P(x' \vert x,a,\rho^{E}_{n-1},\varepsilon^0_{n+1}}_{=1})
        \right) \\[6pt]
        &\le 
        \EE\!\left(
        \sum_{x,a}
        \left\lvert 
        \rho^{E}_{n-1}(x)\bigl(\pi^E_{n-1}(a \vert x,\rho^{E}_{n-1}) - \pi^A_{n-1}(a \vert x,\rho^{E}_{n-1})\bigr)
        +
        (\rho^{E}_{n-1}(x)-\rho^{EA}_{n-1}(x))\,\pi^A_{n-1}(a \vert x,\rho^{E}_{n-1})
        \right\rvert
        \right) \\[6pt]
        &\le
        \EE\!\left(
        \sum_{x,a}
        \rho^{E}_{n-1}(x)\,
        \lvert \pi^E_{n-1}(a \vert x,\rho^{E}_{n-1}) - \pi^A_{n-1}(a \vert x,\rho^{E}_{n-1}) \rvert
        \right)
        +
        \EE\!\left(
        \sum_{x,a}
        \lvert \rho^{E}_{n-1}(x) - \rho^{EA}_{n-1}(x) \rvert\,
        \pi^A_{n-1}(a \vert x,\rho^{E}_{n-1})
        \right) \\[6pt]
        &=
        \underbrace{
        \EE\!\left(
        \EE_{x\sim \rho^E_{n-1}}
        \bigl[
        \lvert
        \pi^E_{n-1}(x,\rho^{E}_{n-1})
        -
        \pi^A_{n-1}(x,\rho^{E}_{n-1})
        \rvert
        \bigr]
        \right)
        }_{=\ \delta^{BC}_{n-1}}
        +
        \EE\!\left(
        \sum_{x}
        \lvert \rho^{E}_{n-1}(x) - \rho^{EA}_{n-1}(x) \rvert\,
        \underbrace{\sum_{a}\pi^A_{n-1}(a \vert x,\rho^{E}_{n-1})}_{=1}
        \right) \\[6pt]
        &\le 
        \delta^{BC}_{n-1} + 
        \EE\!\left( \| \rho^{E}_{n-1} - \rho^{EA}_{n-1} \|_1 \right) \\[4pt]
        &\le 
        n\, \delta^{BC}.
    \end{align*}
    That gives Eq.~\eqref{eq:rhoE-rhoEA_BC}. Now,
    \begin{align*}
        \EE\!\left(\| \mu^{E}_n - \mu^{EA}_n \|_1\right)
        &= 
        \EE\!\left( \sum_{x,a} \lvert \mu^{E}_n(x,a) - \mu^{EA}_n(x,a) \rvert \right) \\[4pt]
        &=
        \EE\!\left(
        \sum_{x,a} 
        \lvert 
        \rho^{E}_{n}(x)\,\pi^E_{n}(a \vert x,\rho^{E}_{n-1})
        -
        \rho^{EA}_{n}(x)\,\pi^A_{n}(a \vert x,\rho^{E}_{n-1})
        \rvert
        \right) \\[4pt]
        &=
        \EE\!\left(
        \sum_{x,a} 
        \lvert 
        \rho^{E}_{n}(x)\bigl(\pi^E_{n}(a \vert x,\rho^{E}_{n-1}) - \pi^A_{n}(a \vert x,\rho^{E}_{n-1})\bigr)
        +
        \pi^A_{n}(a \vert x,\rho^{E}_{n-1})\bigl(\rho^{E}_{n}(x)-\rho^{EA}_{n}(x)\,\bigl)\,
        \rvert
        \right) \\[4pt]
        &\le
        \EE\!\left(
        \sum_{x,a} 
        \rho^{E}_{n}(x)\lvert \pi^E_{n}(a \vert x,\rho^{E}_{n-1}) - \pi^A_{n}(a \vert x,\rho^{E}_{n-1})\rvert \!\right)
        +
        \EE\!\left(\sum_{x,a} 
        \pi^A_{n}(a \vert x,\rho^{E}_{n-1}) \lvert\, \rho^{E}_{n}(x)-\rho^{EA}_{n}(x)\,\rvert
        \!\right) \\[4pt]
        &\le 
        \delta^{BC}_{n-1} + 
        \underbrace{\EE\!\left( \| \rho^{E}_{n} - \rho^{EA}_{n} \|_1 \right)}_{\le n \, \delta^{BC}} \\[6pt]
        &\le 
        (n+1) \, \delta^{BC}.
    \end{align*}
    That gives Eq.~\eqref{eq:muE-muEA_BC}.
\end{proof}
We extend \citep[Lemma 1 C.1]{ramponi2023imitation}. 
\begin{lemma}
    \label{lemma:Vpi1piA-VpiApiA}
    Under Assumption~\eqref{assumption:Lipschitz}~\eqref{assumption:rmax}, for any policies $\pi^A, \pi^1 \in \Pi$, we have 
    \begin{equation}
    \label{eq:Vpi1piA-VpiApiA}
    \begin{split}
        V(\pi^1, \pi^A) - V(\pi^A, \pi^A)\leq&  2 L_r \, \EE\big(\sum_{n = 0}^{H-1} \|\rho^A_n - \rho^E_n\|_1) \\
        &
        +r_{\max} \big[\EE( \sum_{n=0}^{H-1} \| \mu^{E}_n - \mu^{E,A}_n\|\big)  + \EE (\sum_{n= 0}^{H-1} \| \rho^{A,1}_n - \rho^{E,1}_n\|_1) +  \EE (\sum_{n= 0}^{H-1} \| \rho^{A, A}_n - \rho^{E, A}_n\|_1) \\
        &+ \sum_{n= 0}^{H-1}\EE \!\left(\EE_{X \sim \rho^{E, 1}_n}(\|\pi^1(X, \rho^A_n) - \pi^1(X, \rho^E_n))\!\right)   
        + \sum_{n= 0}^{H-1}\EE \!\left(\EE_{X \sim \rho^{E, A}_n}(\|\pi^A(X, \rho^A_n) - \pi^A(X, \rho^E_n)\|)\!\right)\big].
    \end{split}
    \end{equation}
    And by using the Lipschitz property of our policies:
    \begin{equation}
        \label{eq:Vpi1piA-VpiApiA-Lipschitz}
         \begin{split}
            V(\pi^1, \pi^A) - V(\pi^A, \pi^A)\leq&  2 (L_r + L_E \,r_{\max}) \, \EE\big(\sum_{n = 0}^{H-1} \|\rho^A_n - \rho^E_n\|_1) \\
            &
            +r_{\max} \big[\EE( \sum_{n=0}^{H-1} \| \mu^{E}_n - \mu^{E,A}_n\|_1\big)  + \EE (\sum_{n= 0}^{H-1} \| \rho^{A,1}_n - \rho^{E,1}_n\|_1) +  \EE (\sum_{n= 0}^{H-1} \| \rho^{A, A}_n - \rho^{E, A}_n\|_1)
         \end{split}
    \end{equation}
\end{lemma}
\begin{proof}
    Let $\pi^1\in \Pi$, we can write
    \begin{flalign*}
        &V(\pi^1, \pi^A) - V(\pi^A, \pi^A) = \underbrace{V(\pi^1, \pi^A) - V(\pi^E, \pi^E)}_{A} + \underbrace{V(\pi^E, \pi^E) - V(\pi^A, \pi^A) }_{B}.&
    \end{flalign*}
    First, let us work on $A$:
    \begin{flalign*}
        &A = \underbrace{V(\pi^1, \pi^A) - V(\pi^1, \pi^E)}_{A_1} + \underbrace{V(\pi^1, \pi^E) -V(\pi^E, \pi^E) }_{A_2}.&
    \end{flalign*}
    $A_2$ is non positive since $\pi^E$ is an equilibrium. We bound $A_1$ with Lemma~\ref{lemma:Vpi3pi1-Vpi3pi2} and obtain 
    \begin{flalign*}
        &A \leq  L_r \, \EE\big(\sum_{n = 0}^{H-1} \|\rho^A_n - \rho^E_n\|_1) + r_{\max} \EE (\sum_{n= 0}^{H-1} \| \rho^{A,1}_n - \rho^{E,1}_n\|_1)
        +r_{\max}  \sum_{n= 0}^{H-1}\EE (\EE_{X \sim \rho^{E, 1}_n}(\|\pi^1_n(X, \rho^A_n) - \pi^1_n(X, \rho^E_n))). &
    \end{flalign*}
    Consider $B$:
    \begin{flalign*}
        &B \leq \underbrace{V(\pi^E, \pi^E) - V(\pi^A, \pi^E)}_{B_1} + \underbrace{V(\pi^A, \pi^E) - V(\pi^A, \pi^A)}_{B_2} &
    \end{flalign*}
    \begin{flalign*}
        B_1 &\leq \EE\left(\sum_{n=0}^{H-1} \sum_{x,a}\big( \mu^{E, E}_n(x,a) r(x, a, \rho^E) - \mu^{E,A}_n(x,a)r(x, a, \rho^E)\big) \right) &\\
        &\leq r_{\max}\EE\left( \sum_{n=0}^{H-1} \sum_{x,a}\big| \mu^{E, E}_n(x,a) - \mu^{E,A}_n(x,a) \big|\right)& \\
        &\leq r_{\max} \EE( \sum_{n=0}^{H-1} \| \mu^{E, E}_n - \mu^{E,A}_n\|\big).&
    \end{flalign*}
    By Lemma~\ref{lemma:Vpi3pi1-Vpi3pi2} we obtain 
    \begin{equation*}
        B_2 \leq L_r \, \EE\big(\sum_{n = 0}^{H-1} \|\rho^A_n - \rho^E_n\|_1) + r_{\max} \EE (\sum_{n= 0}^{H-1} \| \rho^{A, A}_n - \rho^{E, A}_n\|_1)
        +r_{\max}  \sum_{n= 0}^{H-1}\EE (\EE_{x \sim \rho^{E, A}}(\|\pi^A_n(x, \rho^A_n) - \pi^A_n(x, \rho^E_n))). 
    \end{equation*}
    Summing everything we obtain Eq.~\eqref{eq:Vpi1piA-VpiApiA}.
\end{proof}

We give two new lemmas. 
\begin{lemma}
    \label{lemma:rhoA-rhoE_LP}
    \begin{equation}
        \label{eq:rhoA-rhoE_BC_LP=0}
        \EE(\sum_{n= 0}^{H-1}\|\rho^A_n - \rho^E_n \|_1) \leq \delta^{BC}\frac{(1 + L_E + L_P)^H}{ (L_E + L_P)^2},
    \end{equation}
    \begin{equation}
        \label{eq:rhoA-rhoE_ADV_LP}
        \EE(\|\rho^A_n - \rho^E_n \|_1) \leq \EE(\|\mu_{n}^A - \mu_{n}^E\|_1) \leq \delta^{ADV} .
    \end{equation}
\end{lemma}
\begin{proof}
    \begin{align*}
        \EE(\|\rho^A_{n+1}-\rho^E_{n+1}\|) &= \EE(\sum_{x', x, a}\lvert \pi^A_{n}(a\vert x, \rho_n^{A}) \rho^A_n(x)  P(x'\vert x, a,\rho^A_n,  \varepsilon^0_{n+1}) - \pi^E_{n}(a\vert x, \rho_n^{E}) \rho^E_n(x)  P(x'\vert x, a,\rho^E_n \varepsilon^0_{n+1})\rvert ) \\
        &\leq \EE(\sum_{x', x, a}\pi^A_{n}(a\vert x, \rho_n^{A}) \rho^A_n(x)  \lvert P(x'\vert x, a,\rho^A_n,  \varepsilon^0_{n+1}) - P(x'\vert x, a,\rho^E_n,  \varepsilon^0_{n+1}) \rvert  ) \\
        &+ \EE(\sum_{x',x, a}P(x'\vert x, a,\rho^E_n,  \varepsilon^0_{n+1}) \lvert \pi^A_{n}(a\vert x, \rho_n^{A}) \rho^A_n(x) - \pi^E_{n}(a\vert x, \rho_n^{E}) \rho^E_n(x) \rvert ) \\
        &\leq \underbrace{\EE(\sum_{x, a} \pi^A_{n}(a\vert x, \rho_n^{A}) \rho^A_n(x)  \| P( x, a,\rho^A_n,  \varepsilon^0_{n+1}) - P( x, a,\rho^E_n,  \varepsilon^0_{n+1}) \|_1 )}_{\EE(\EE_{(x, a)\sim\mu_n^A}(\| P( x, a,\rho^A_n,  \varepsilon^0_{n+1}) - P( x, a,\rho^E_n,  \varepsilon^0_{n+1}) \|_1))} \\
        &+ \EE(\sum_{x, a}\lvert \pi^A_{n}(a\vert x, \rho_n^{A}) \rho^A_n(x) - \pi^E_{n}(a\vert x, \rho_n^{E}) \rho^E_n(x) \rvert ) \\
        &\leq L_P \EE(\|\rho^A_n - \rho^E_n\|_1) + \underbrace{\EE(\sum_{x, a}\lvert \pi^A_{n}(a\vert x, \rho_n^{A}) \rho^A_n(x) - \pi^A_{n}(a\vert x, \rho_n^{A}) \rho^E_n(x) \rvert )}_{A_1} \\
        &+ \underbrace{\EE(\sum_{x, a}\lvert \pi^A_{n}(a\vert x, \rho_n^{A}) \rho^E_n(x) - \pi^A_{n}(a\vert x, \rho_n^{E}) \rho^E_n(x) \rvert )}_{A_2} \\
        &+ \underbrace{\EE(\sum_{x, a}\lvert \pi^A_{n}(a\vert x, \rho_n^{E}) \rho^E_n(x) - \pi^E_{n}(a\vert x, \rho_n^{E}) \rho^E_n(x) \rvert )}_{A_3}
    \end{align*}
    First we have $A_1 = \EE(\|\rho^A_n - \rho^E_n\|_1).$ Then we have 
    \begin{flalign*}
        &A_2 = \EE\big(\EE_{x\sim \rho^E_n}(\|\pi^A_n(x, \rho^A_n) - \pi^A_n(x, \rho^E_n)\|_1)\big) \leq L_E \,\EE(\|\rho^A_n - \rho^E_n\|_1).&
    \end{flalign*}
    Now,
    \begin{flalign*}
        &A_3 = \underbrace{\EE\big(\EE_{x\sim\rho_n^E}( \| \pi^A_n( x, \rho^E_n) - \pi^E_n(x, \rho^E_n)\|\big)}_{\leq \delta^{BC}}. &\\
    \end{flalign*}
    By combining the terms we have 
    \begin{equation}
        \EE(\|\rho^A_{n+1}-\rho^E_{n+1}\|) \leq \EE(\|\rho^A_{n}-\rho^E_{n}\|)(1 + L_E + L_P) + \delta^{BC} \leq  \delta^{BC} \sum^{n}_{k = 0}(1+L_E + L_P)^k = \delta^{BC} \frac{(1 + L_E + L_P)^n - 1}{L_E + L_P}. 
    \end{equation}
    By summing over $n$ we obtain the result. For Eq.~\eqref{eq:rhoA-rhoE_ADV_LP}, we have 
    \begin{flalign*}
        \EE(\|\rho^A_n - \rho^E_n\|_1) &\leq \EE(\sum_{x}\rvert\rho_n^A(x) - \rho^E_n(x)\lvert ) &\\
        &\leq \EE(\sum_x\rvert\sum_a\mu^A_n(x, a) - \mu^E_n(x, a)\lvert) \\
        &\leq \EE(\sum_{x, a}\rvert \mu^A_n(x, a) - \mu^E_n(x, a)\lvert) \\
        &\leq \EE(\|\mu_n^A - \mu_n^E\|_1) \leq \delta^{ADV}.
    \end{flalign*}
\end{proof}
\begin{lemma}
    \label{lemma:rhoA1-rhoE1_LP=0}
    Let $\pi^1 \in \Pi$. We have 
    \begin{equation}
        \label{eq:rhoA1-rhoE1_BC_LP}
        \EE(\sum_{n=0}^{H-1} \| \rho_n^{A, 1} - \rho_n^{E, 1}\|) \leq \delta^{BC} \frac{ (1 + L_E + L_P)^{H}}{(L_P + L_E)^2}.
    \end{equation}
    and 
    \begin{equation}
        \label{eq:rhoA1-rhoE1_ADV_LP}
        \EE(\sum_{n=0}^{H-1} \| \rho_n^{A, 1} - \rho_n^{E, 1}\|) \leq  \delta^{ADV}(L_P + L_E)H^2.
    \end{equation}
\end{lemma}
\begin{proof}
    Suppose $L_P + L_E \ne 0$,
    \begin{align*}
        \EE(\| \rho_n^{A, 1} - \rho_n^{E, 1}\|) &= \EE \!\left(\sum_{x'} \vert\rho_n^{A, 1}(x') - \rho_n^{E, 1}(x')\vert\!\right) \\
        &\leq \EE \!\left(\sum_{x', x, a} \vert \rho_{n-1}^{A, 1}(x) \pi^1_{n-1}(a\vert x, \rho^A_{n-1})P(x'\vert x, a,\rho^A_{n-1} \varepsilon^0_{n})  - \rho_{n-1}^{E, 1}(x)\pi^1_{n-1}(a\vert x, \rho^E_{n-1})P(x'\vert x, a,\rho^E_{n-1} \varepsilon^0_{n}) \vert\!\right) \\
        &\leq \EE \!\left(\sum_{x', x, a} \rho_{n-1}^{A, 1}(x) \pi^1_{n-1}(a\vert x, \rho^A_{n-1})\vert P(x'\vert x, a,\rho^A_{n-1}, \varepsilon^0_{n})  - P(x'\vert x, a,\rho^E_{n-1}, \varepsilon^0_{n}) \vert\!\right) \\
        & + \EE \!\left(\sum_{x', x, a}  P(x'\vert x, a, \rho^E_{n-1} ,\varepsilon^0_{n+1}) \vert\rho_{n-1}^{A, 1}(x) \pi^1_{n-1}(a\vert x, \rho^A_{n-1}) - \rho_{n-1}^{E, 1}(x)\pi^1_{n-1}(a\vert x, \rho^E_{n-1})\vert\!\right) \\
        &\leq \underbrace{\EE \!\left(\sum_{x, a} \rho_{n-1}^{A, 1}(x) \pi^1_{n-1}(a\vert x, \rho^A_{n-1})\| P(x, a,\rho^A_{n-1}, \varepsilon^0_{n})  - P(x, a,\rho^E_{n-1}, \varepsilon^0_{n})\|_1\right)}_{A_1} \\
        & + \underbrace{\EE \!\left(\sum_{x, a}  \vert\rho_{n-1}^{A, 1}(x) \pi^1_{n-1}(a\vert x, \rho^A_{n-1}) - \rho_{n-1}^{E, 1}(x)\pi^1_{n-1}(a\vert x, \rho^E_{n-1})\vert\!\right)}_{A_2} \\
    \end{align*}
    Let us consider $A_1$
    \begin{flalign*}
        &A_1 = \EE\!\left( \EE_{(x, a)\sim\mu^{A, 1}_{n-1}}(\| P(x, a,\rho^A_{n-1}, \varepsilon^0_{n})  - P(x, a,\rho^E_{n-1}, \varepsilon^0_{n})\|_1) \!\right) \leq L_P\EE(\| \rho_{n-1}^{A} - \rho_{n-1}^{E}\|). &
    \end{flalign*}
    Let us consider $A_2$
    \begin{flalign*}
        A_2 &=\EE \!\left(\sum_{x, a}  \vert\rho_{n-1}^{A, 1}(x) \pi^1_{n-1}(a\vert x, \rho^A_{n-1}) - \rho_{n-1}^{E, 1}(x)\pi^1_{n-1}(a\vert x, \rho^E_{n-1})\vert\!\right) &\\
        &\leq \EE \!\left(\sum_{x, a}  \vert\pi^1_{n-1}(a\vert x, \rho^A_{n-1})\big(\rho_{n-1}^{A, 1}(x) - \rho_{n-1}^{E, 1}(x) \big)  + \rho_{n-1}^{E, 1}(x)\big(\pi^1_{n-1}(a\vert x, \rho^E_{n-1}) - \pi^1_{n-1}(a\vert x, \rho^A_{n-1})\big)\vert\!\right) \\
        &\leq \EE \!\left(\sum_{x, a}  \pi^1_{n-1}(a\vert x, \rho^A_{n-1})\vert\rho_{n-1}^{A, 1}(x) - \rho_{n-1}^{E, 1}(x) \vert \!\right) + \EE\!\left(\sum_{x,a}\rho_{n-1}^{E, 1}(x)\vert\pi^1_{n-1}(a\vert x, \rho^E_{n-1}) - \pi^1_{n-1}(a\vert x, \rho^A_{n-1})\vert\!\right) \\
        &\leq \EE(\| \rho_{n-1}^{A, 1} - \rho_{n-1}^{E, 1}\|) + \EE(\EE_{x \sim \rho^{E, 1}}(\|\pi^1_{n-1}(x, \rho^E_{n-1}) - \pi^1_{n-1}(x, \rho^A_{n-1}) \|_1) \\
        &\leq \EE(\| \rho_{n-1}^{A, 1} - \rho_{n-1}^{E, 1}\|) + L_E \EE(\|\rho^E_{n-1} - \rho^A_{n-1} \|_1) \\
    \end{flalign*}
    Combining $A_1$ and $A_2$ we obtain 
    \begin{align*}
         \EE(\| \rho_n^{A, 1} - \rho_n^{E, 1}\|) &\leq \EE(\| \rho_{n-1}^{A, 1} - \rho_{n-1}^{E, 1}\|)  + (L_E + L_P) \EE(\|\rho^E_{n-1} - \rho^A_{n-1} \|_1) \\
         &\leq  \EE(\| \rho_{n-1}^{A, 1} - \rho_{n-1}^{E, 1}\|) +  \delta^{BC} (L_E + L_P)  \frac{(1 + L_E + L_P)^{n-1} - 1}{L_E + L_P} \\
         &\leq  \EE(\| \rho_{n-1}^{A, 1} - \rho_{n-1}^{E, 1}\|) +  \delta^{BC} (1 + L_E + L_P)^{n-1}
    \end{align*}
    By summing we obtain 
    \begin{equation*}
        \EE(\| \rho_n^{A, 1} - \rho_n^{E, 1}\|) \leq \delta^{BC} \frac{ (1 + L_E + L_P)^{n}}{L_P + L_E}
    \end{equation*}
    and by summing again
    \begin{equation*}
        \EE(\sum_{n=0}^{H-1} \| \rho_n^{A, 1} - \rho_n^{E, 1}\|) \leq\delta^{BC} \frac{ (1 + L_E + L_P)^{n}}{(L_P + L_E)^2}.
    \end{equation*}
    Let us prove Eq.~\eqref{eq:rhoA1-rhoE1_ADV_LP}. Previously, one of the intermediate step showed that $\EE(\| \rho_n^{A, 1} - \rho_n^{E, 1}\|) \leq  \EE(\| \rho_{n-1}^{A, 1} - \rho_{n-1}^{E, 1}\|) +(L_P + L_E) \EE(\|\rho^E_{n-1} - \rho^A_{n-1} \|_1)$. Applying Lemma~\ref{lemma:rhoA-rhoE_LP} with Eq.~\eqref{eq:rhoA-rhoE_ADV_LP} leads to 
    \begin{equation*}
        \EE(\| \rho_n^{A, 1} - \rho_n^{E, 1}\|) \leq \EE(\| \rho_{n-1}^{A, 1} - \rho_{n-1}^{E, 1}\|) + (L_E + L_P)\delta^{ADV},
    \end{equation*}
    from what we deduce the result by summing twice.
\end{proof}

\begin{lemma}
    \label{lemma:muEA-muEE_ADV_LP}
    We have 
    \begin{equation}
        \label{eq:muEA-muEE_ADV_LP}
        \EE(\sum_{n=1}^{H-1} \|\mu_n^{EA} - \mu_n^E \|_1) \leq \delta^{ADV}(H(L_E +1) + H^2(L_P + L_E)).
    \end{equation}
\end{lemma}
\begin{proof}
    \begin{align*}
        \EE(|\mu_n^{EA} - \mu_n^E \|_1) &\leq \EE(\|\mu_n^{EA} - \mu_n^A \|_1) + \underbrace{\EE(\|\mu_n^{A} - \mu_n^E \|_1)}_{\leq\,\delta^{ADV}}.
    \end{align*}
    Now we have 
    \begin{align*}
        \EE(\|\mu_n^{EA} - \mu_n^A \|_1) &\leq  \EE(\sum_{x,a} \lvert\rho^{EA}_{n}(x)\pi^A_n(a\vert x, \rho^E_n) - \rho^{A}_{n}(x)\pi^A_n(a\vert x, \rho^A_n) \rvert) \\
         &\leq  \underbrace{\EE(\sum_{x,a} \rho^{EA}_{n}(x)\lvert\pi^A_n(a\vert x, \rho^E) - \pi^A_n(a\vert x, \rho^A_n)\vert  )}_{=\, \EE\big(\EE_{X\sim\rho_n^{EA}}(\|\pi^A(X, \rho_n^E) - \pi^A(X, \rho_n^A)\|)\big)}+ \EE(\sum_{x,a} \pi^A_n(a\vert x, \rho^A_n)\vert\rho^{EA}_{n}(x) - \rho^{A}_{n}(x)\vert) \\
         &\leq L_E\EE(\|\rho^E_n - \rho^A_n\|_1) + \EE(\|\rho^{EA}_n - \rho^A_n\|_1).
    \end{align*}
    We deduce that 
    \begin{align*}
        \EE(|\mu_n^{EA} - \mu_n^E \|_1) &\leq L_E\EE(\|\rho^E_n - \rho^A_n\|_1) + \EE(\|\rho^{EA}_n - \rho^A_n\|_1) + \delta^{ADV}.
    \end{align*}
    Applying Lemma~\ref{lemma:rhoA-rhoE_LP} with Eq.~\eqref{eq:rhoA-rhoE_ADV_LP} and Lemma~\ref{lemma:rhoA1-rhoE1_LP=0} with Eq.~\eqref{eq:rhoA1-rhoE1_ADV_LP} we obtain
    \begin{align*}
        \EE(\sum_{n=1}^{H-1} \|\mu_n^{EA} - \mu_n^E \|_1) \leq  \delta^{ADV}(H(L_E +1) + H^2(L_P + L_E)),
    \end{align*}
\end{proof}

\subsection{Proofs of the Theorems}

To prove the theorems, it is enough to apply the previous established inequalities. 

\noindent\textbf{Proof of Theorem~\ref{theorem:BC_LPge0}}

\begin{proof}
    Apply Lemma~\ref{lemma:Vpi1piA-VpiApiA} with Eq.~\eqref{eq:Vpi1piA-VpiApiA-Lipschitz}, then Lemma~\ref{lemma:rhoE-rhoEA_BC-muE-muEA_BC}, Lemma~\ref{lemma:rhoA-rhoE_LP} with Eq.~\eqref{eq:rhoA-rhoE_BC_LP=0} and Lemma~\ref{lemma:rhoA1-rhoE1_LP=0} with Eq.~\eqref{eq:rhoA1-rhoE1_BC_LP}.
\end{proof}

\noindent\textbf{Proof of Theorem~\ref{theorem:ADV_LPge0}}

\begin{proof}
    Apply Lemma~\ref{lemma:Vpi1piA-VpiApiA} with Eq.~\eqref{eq:Vpi1piA-VpiApiA-Lipschitz}, then Lemma~\ref{lemma:muEA-muEE_ADV_LP}, Lemma~\ref{lemma:rhoA-rhoE_LP} with Eq.~\eqref{eq:rhoA-rhoE_ADV_LP} and Lemma~\ref{lemma:rhoA1-rhoE1_LP=0} with Eq.~\eqref{eq:rhoA1-rhoE1_ADV_LP}.
\end{proof}

\noindent\textbf{Proof of Theorem~\ref{theorem:BC_LPge0_reward}}

\begin{proof}
    In the proof of Lemma~\ref{lemma:Vpi1piA-VpiApiA}, we show that 
    % \vspace{-10pt}
    $$
    \vert V(\pi^E, \pi^E) - V(\pi^A, \pi^E) \vert \leq r_{\max}\EE\left[ \sum_{n=0}^{H-1} \| \mu^E_n - \mu^{E, A} \|_1 \right].
    % \vspace{-8pt}
    $$
    We can conclude by applying Lemma~\ref{lemma:rhoE-rhoEA_BC-muE-muEA_BC}.
\end{proof}

\noindent\textbf{Proof of Theorem~\ref{theorem:ADV_LPge0_reward}}

\begin{proof}
    The proof uses the same argument as in  Theorem~\ref{theorem:BC_LPge0_reward}, but by using Lemma~\ref{lemma:muEA-muEE_ADV_LP}.
\end{proof}



\clearpage

\section{Additional Algorithms}
\label{appendix:algo}

\begin{algorithm}[hbp]
\caption{Proxy Algorithm}
\begin{algorithmic}[1]
\label{algo:proxy}
\REQUIRE MFG $\cM = (\cX, \cA, \cE^0, \nu^0, P, r, H, \rho_0)$; \\ 
\hspace{3.5em} $\pi^A, \pi^E$; \\
\hspace{3.5em} Monte Carlo parameter $N$;
\STATE Initialize the empirical BC proxy for $t = 0, \dots, H-1$, $\hat \delta^{BC}_t = 0$.
\STATE Initialize the empirical ADV proxy for $t = 0, \dots, H-1$, $\hat \delta^{ADV}_t = 0$.
\STATE Initialize $$\hat \delta^{BC}_0 = \sum_{x} \rho_0(x) \|\pi^A_0(x, \rho_0) - \pi^E_0(x, \rho_0)\|_1.$$
\STATE Initialize $$\hat \delta^{ADV}_0 =  \sum_{x, a} \rvert\rho_0(x)\pi^A_0(a\vert x, \rho_0) - \rho_0(x)\pi^E_0(a\vert x, \rho_0)\lvert.$$
\FOR{$n = 1,\dots,N$}
    \STATE Sample one i.i.d sequence of common noises $\varepsilon^{0, n}_1, \dots, \varepsilon^{0, n}_{H-1}$ with $\varepsilon^{0, n}_1 \sim \nu^0$.
    \STATE Initialize the mean-field $\rho_0^{E,n} = \rho_0$.
    \STATE Initialize the mean-field $\rho_0^{A,n} = \rho_0$.
    \FOR{$t = 1,\dots,H-1$}
        \STATE Compute $$\rho_t^{E,n} = \sum_{x, a} \rho_{t-1}^{E,n}(x)\pi^E_{t-1}(a\mid x, \rho^{E,n}_{t-1}) P(\cdot\mid x, a, \rho^{E,n}_{t-1}, \varepsilon^{0,n}_t).$$
        \STATE Compute $$\rho_t^{A,n} = \sum_{x, a} \rho_{t-1}^{A,n}(x)\pi^{A,n}_{t-1}(a\mid x, \rho^{A,n}_{t-1}) P(\cdot\mid x, a, \rho^{A,n}_{t-1}, \varepsilon^{0,n}_t).$$
        \STATE Update $$\hat \delta^{BC}_t += \frac{1}{N} \sum_{x} \rho_t^{E,n}(x) \|\pi^A_t(x, \rho_t^{E,n}) - \pi^E_t(x, \rho_t^{E,n})\|_1.$$
        \STATE Update $$\hat \delta^{ADV}_t += \frac{1}{N} \sum_{x, a} \rvert\rho_t^{A,n}(x)\pi^A_t(a\vert x, \rho_t^{A,n}) - \rho_t^{E,n}(x)\pi^E_t(a\vert x, \rho_t^{E,n})\lvert.$$
    \ENDFOR
\ENDFOR
\STATE Compute $\hat \delta^{BC} = \max_{t \in [H-1]} \hat \delta^{BC}_t$.
\STATE Compute $\hat \delta^{ADV} = \max_{t \in [H-1]} \hat \delta^{ADV}_t$.
\STATE \textbf{Return} $\hat \delta^{BC}, \hat \delta^{ADV}$
\end{algorithmic}
\end{algorithm}

\begin{algorithm}[htbp]
\caption{Neural Network Best Response}
\begin{algorithmic}[1]
\label{algo:nn_best_response}
\REQUIRE MFG $\cM = (\cX, \cA, \cE^0, \nu^0, P, r, H, \rho_0)$; Sequence of policies $(\pi^0, \dots, \pi^k) \in \Pi^{k+1}$; Iterations $J$; Batch size $B$; Optimizer. \\ 
\FOR{$j = 1, \dots, J$} 
    \STATE Reset batch loss: $L(\theta_{j-1}) = 0$.
    \STATE Sample a batch of i.i.d. common noises $(\varepsilon^{0, s}_t)_{0 \leq t \leq H-1}$ for $s=1, \dots, B$.
    \FOR{$s = 1, \dots, B$}
        \STATE Compute the aggregated sequence $\bar{\rho}^{k,s}$ induced by $(\pi^0, \dots, \pi^k)$.
        \STATE Compute the mean-field sequence $\rho^{\theta, s}$ induced by the learner $\pi_\theta$ responding to $\bar{\rho}^{k,s}$.
        \STATE Accumulate the loss:
        \[
        L(\theta_{j-1}) \leftarrow L(\theta_{j-1}) - \frac{1}{B} \sum_{t=0}^{H-1} \sum_{x, a} \rho^{\theta,s}_t(x) \pi^\theta_t(a \mid x, \bar{\rho}^{k,s}_t) r(x, a, \bar{\rho}^{k,s}_t)
        \]
    \ENDFOR
    \STATE Compute the gradient $\nabla_\theta L(\theta_{j-1})$ and update parameters via the optimizer to obtain $\theta_{j}$. 
\ENDFOR 
\STATE \textbf{Return} Optimized parameters $\theta_{J}$ defining the best response policy $\pi_{\theta_J}$.
\end{algorithmic}
\end{algorithm}

%%%%%%%%%%%%%%%%%%%%%%%%%%%%%%%%%%%%%%%%%%%%%%%%%%%%%%%%%%%%%%%%%%%%%%%%%%%%%%%
%%%%%%%%%%%%%%%%%%%%%%%%%%%%%%%%%%%%%%%%%%%%%%%%%%%%%%%%%%%%%%%%%%%%%%%%%%%%%%%

\clearpage

\section{Additional Details to Example 1}
\label{appendix:example1}

\subsection{Mann Iterations and Best Response}
We use the following algorithm to compute a $\epsilon$-Nash equilibrium to obtain the expert policy $\pi^E$. 
\begin{algorithm}[ht]
\caption{Mann Iteration for $\epsilon$-Nash Equilibrium}
\begin{algorithmic}[1]
\label{algo:mann_iteration_mfg}
\REQUIRE MFG environment $\cM$; Initial policy $\pi^0 \in \Pi$, Learning rate sequence $(\gamma_k)_{1\leq k\leq K} \in (0, 1]$; Number of iterations $K$. \\ 
\STATE Define $\pi^{0,*} = \pi^0$
\FOR{each $k = 0, \dots, K-1$}
    \STATE Compute best response $\pi^{BR, k}$ to $\pi^{k-1,*}$. 
    \STATE Define $\pi^{k,*} = \pi^{k-1,*} + \gamma_k \pi^{BR, k}$
\ENDFOR
\STATE \textbf{Return} $\pi^{*,K}$.
\end{algorithmic}
\end{algorithm}

To compute a best response, we discretize the mean-field space, and use a Monte Carlo estimation to perform a backward induction. 

\begin{algorithm}[ht]
\caption{Backward Induction for Best Response.}
\begin{algorithmic}[1]
\label{algo:best_response_dp}
\REQUIRE Fixed population policy $\pi$; MFG $\cM$; Grid for mean-field $\mathcal{G}_{\rho}$; Monte Carlo batch size $M$.
\STATE \textbf{Initialize} Value function $V_H(x, \rho) = 0$ for all $x \in \cX, \rho \in \mathcal{G}_{\rho}$.
\FOR{time $t = H-1, \dots, 0$}
    \FOR{each possible density $\rho \in \mathcal{G}_{\rho}$}
        \FOR{each state $x \in \cX$}
            \FOR{each action $a \in \cA$}
                \STATE Sample a batch of common noise $(\varepsilon^{0, s})_{1\leq s\leq M}$.
                \FOR{each $s \in [M]^*$}
                    \STATE Compute expected next density: $\rho^{'s} = \Phi(\rho^s, \pi_t, \varepsilon^{0,s}_{t+1})$.
                \ENDFOR
                \STATE Define $Q_t(x, \rho, a) = r(x, a, \rho) + \sum_{s = 1}^M  V_{t+1}(x', \rho^{'s})\rho^{'s}(x')$
            \ENDFOR
            \STATE $V_t(x, \rho) = \max_{a \in \cA} Q_t(x, \rho, a)$
            \STATE $\pi^{BR}_t(a \mid x, \rho) = \mathds{1}_{\{a = \text{argmax}_{a'} Q_t(x, \rho, a')\}}$
        \ENDFOR
    \ENDFOR
\ENDFOR
\STATE \textbf{Return} $\pi^{BR} = (\pi^{BR}_t)_{0 \leq t < H}$.
\end{algorithmic}
\end{algorithm}

\subsection{Vanilla and Adaptive Policies}

We describe here how we obtain using data from expert agents following the policy $\pi^E$, the vanilla and adaptive policies used for the experiments in Subsection~\ref{subsection:example1}.

\noindent\textbf{Data Generation:}
The simplicity of the environment allows us to consider classical statistical estimators for the vanilla and adaptive policies. Let $(\eta, \alpha)$ be a couple of parameters. We generate a synthetic dataset composed of $N$ independent trajectories of $M$ agents:
$
\cD_{\alpha,\eta}^{N, M} = \left\{ \left( x^{n,m}_{t, \alpha, \eta}, a^{n,m}_{t, \alpha, \eta} \right)_{\substack{0\leq t\leq H-1 \\ 0\leq m \leq M-1}} \right\}_{1\leq n\leq N},
$
where each agent acts according to the expert policy $\pi^E_{\alpha, \eta}$ in an environment subjected to a common noise realization $\varepsilon^{0,n}_{1:H-1}$. Specifically, for each trajectory $n$, a unique common noise sequence $(\varepsilon^{0,n}_{t})_{0 \leq t \leq H-1}$ is sampled, inducing a specific mean-field sequence $\rho^{E,n}_{\alpha, \eta}(\varepsilon^{0,n}_{1:H-1})$ that the agents observe and react to. To simplify notation, we fix the parameters $(\alpha, \eta)$ and suppress their explicit dependence in the following subscripts.

\noindent\textbf{Vanilla Policy:}  
We define the vanilla policy $\hat \pi^{M, N,\text{vanilla}}$ as the time-dependent approximation that ignores the common noise fluctuations. It is computed via the empirical marginals:
% \vspace{-14pt}
$$
\hat\pi_t^{M, N,\text{vanilla}}(a\mid x) = \frac{\hat \mu^{E, M, N}_t(x, a)}{\sum_{a' \in \mathcal{A}} \hat \mu^{E, M, N}_t(x, a')},
% \vspace{-12pt}
$$

where $\hat \mu^{E, M, N}_t(x, a) = \frac{1}{NM}\sum_{n,m} \mathds{1}_{\{x^{n,m}_t = x, a^{n,m}_t = a\}}$.

\noindent\textbf{Adaptive Policy:} 
We define the adaptive policy $\hat\pi^{M, N, \text{adaptive}}$ as a density-dependent policy. To reconstruct the policy for any arbitrary mean-field value $\rho$, we employ a Nadaraya–Watson kernel regression:
% \vspace{-7pt}
\begin{equation*}
\begin{split}
        &\hat\pi_t^{M, N, \text{adaptive}}(a\mid x, \rho) = \\ &\frac{\sum_{n, m}\mathds{1}_{\{x^{n,m}_t = x, a^{n,m}_t = a\}} K_h(\hat \rho^{E, n, M}_t - \rho)}{\sum_{a' \in \mathcal{A}}\sum_{n, m}\mathds{1}_{\{x^{n,m}_t = x, a^{n,m}_t = a'\}} K_h(\hat  \rho^{E, n, M}_t  - \rho)},
\end{split}
\end{equation*}
where $ \hat \rho^{E, n, M}_t(1) = \frac{1}{M}\sum_{m} \mathds{1}_{\{x^{n,m}_t = 1\}}$ is the empirical density of agents in state 1 for trajectory $n$. The kernel function $K_h:\Delta_\cX\times\Delta_\cX\to\RR_+$ is defined by the Gaussian radial basis function:
% \vspace{-10pt}
$$
K_h(\rho, \rho') = \exp\left(-\frac{\|\rho - \rho'\|^2_2}{2h^2}\right), \quad \forall (\rho, \rho') \in \Delta_\cX^2. 
% \vspace{-9pt}
$$
This kernel density estimation (KDE) approach allows the agent to interpolate the expert's behavior across the continuous space of possible mean-field realizations.

\subsection{Hyperparameters}

All the experiments have been run with parameters presented in Table~\ref{tab:ex1_hyperparameters}. 

\begin{table}[htbp]
\centering
\caption{Summary of Hyperparameters for the First Example}
\label{tab:ex1_hyperparameters}
\begin{tabular}{@{}llll@{}}
\toprule
\textbf{Category} & \textbf{Parameter} & \textbf{Symbol} & \textbf{Value} \\ \midrule
\textbf{Nash Equilibrium} & Mann Iterations & $K$ & $50$ \\
& Learning Rate & $\gamma_k$ & $0.05$ \\
& Mean Field Discretizations & $\mathcal{G}_{\rho}$ & $50$ points \\ \addlinespace
& MC Samples (Backward Induction) & - & $10000$ \\ \addlinespace
\textbf{Data Generation} & Number of Trajectories & $N$ & $2000$ \\
& Number of Agents & $M$ & $100$ \\ \addlinespace
\textbf{Kernel Regression} & Bandwidth & $h$ & $0.05$ \\ \bottomrule
\end{tabular}
\end{table}

\subsection{Full Sensitivity Analysis }

Figure \ref{fig:ex1_grid_mean} represents the average over $5$ runs of the metrics, as functions of the varying parameters $\eta$ and $\alpha$. Figure \ref{fig:ex1_grid_std} shows the standard deviation. For the row corresponding to the proxy, the advantage corresponds to the adaptive proxy minus the vanilla proxy. For the rows corresponding to the reward vs expert, it corresponds to the reward of the adaptive policy minus the one of the vanilla policy. For the rows corresponding to the exploitability, it corresponds to the exploitability of the vanilla minus the one of the adaptive policy.

% \begin{figure}[htbp]
%     \centering
%     \includegraphics[width=0.9\linewidth]{Figures/Example1/ex1_alpha=1.75_FULL.pdf}
%     \caption{Performance metrics for $5$ different runs, with $\alpha = 1.75$.}
%     \label{fig:example1_alpha=1.75}
% \end{figure}

\begin{figure}[htpb]
    \centering
    % Left plot: Variation with eta
    \begin{subfigure}{0.49\linewidth}
        \centering
        \includegraphics[width=\linewidth]{Figures/Example1/ex1_eta=0.75_FULL.pdf}
        \caption{Fixed $\eta = 0.75$.}
        \label{fig:example1_sub_eta}
    \end{subfigure}
    \hfill
    % Right plot: Variation with alpha
    \begin{subfigure}{0.49\linewidth}
        \centering
        \includegraphics[width=\linewidth]{Figures/Example1/ex1_alpha=1.75_FULL.pdf}
        \caption{Fixed $\alpha = 1.75$.}
        \label{fig:example1_sub_alpha}
    \end{subfigure}

    \caption{Performance metrics for 5 runs in the Example 1. Left: variation across $\alpha$ with fixed $\eta$; Right: variation across $\eta$ with fixed $\alpha$.}
    \label{fig:example1_comparison_side_by_side}
\end{figure}

\begin{figure}[htbp]
    \centering
    \includegraphics[width=0.9\linewidth]{Figures/Example1/Ex1_GRID_FULL.pdf}
    \caption{Example 1: Metrics Averages}
    \label{fig:ex1_grid_mean}
\end{figure}

\begin{figure}[htbp]
    \centering
    \includegraphics[width=0.9\linewidth]{Figures/Example1/Ex1_GRID_STD_FULL.pdf}
    \caption{Example 1: Metrics STD}
    \label{fig:ex1_grid_std}
\end{figure}

\clearpage

\section{Additional Details to the Beach Bar Example}
\label{appendix:beachbar}
\subsection{Hyperparameters}
All the experiments have been run with the parameters presented in Table~\ref{tab:beachbar_hyperparameters}
\begin{table}[htbp]
\centering
\caption{Summary of Hyperparameters for the Beach Bar Environment}
\label{tab:beachbar_hyperparameters}
\begin{tabular}{@{}lll@{}}
\toprule
\textbf{Category} & \textbf{Parameter} & \textbf{Value} \\ \midrule
\textbf{Neural Architecture} & Model Type & MLP \\
& Hidden Layers & $2$ \\
& Layer Width & $64$ \\
& Activation Function & ReLU \\
& Optimizer & Adam \\ \addlinespace
\textbf{Fictitious Play} & Number of (FP) iterations & $40$ \\ \addlinespace
\textbf{Best Response Training} & Number of iterations & $1\,000$ \\
& Learning Rate & $10^{-4}$ \\
& Batch Size  & $500$ \\ \addlinespace
\textbf{Expert Imitation} & Number of iterations & $10\,000$ \\
& Learning Rate & $10^{-4}$ \\
& Batch Size & $200$ \\ \addlinespace
\textbf{Vanilla and Adaptive Learning} & Number of iterations & $10\,000$ \\
& Learning Rate & $10^{-4}$ \\
& Batch Size & $50$ \\ 
& Number of agents in the Simulation & $1000$ \\ 
\bottomrule
\end{tabular}
\end{table}

\subsection{Full Sensitivity Analysis and Training Losses}
Figure \ref{fig:beachbar_grid_mean} represents the average over $5$ runs of the metrics, as functions of the varying parameters $\eta$ and $\alpha$. Figure \ref{fig:beachbar_grid_std} shows the standard deviation. For the row corresponding to the proxy, the advantage corresponds to the adaptive proxy minus the vanilla proxy. For the rows corresponding to the reward vs expert, it corresponds to the reward of the adaptive policy minus the one of the vanilla policy. For the rows corresponding to the exploitability, it corresponds to the exploitability of the vanilla minus the one of the adaptive policy.

Figure \ref{fig:beachbar_losses} illustrates the convergence of Fictitious Play (FP) with common noise (Algorithm \ref{algo:fictitious}), the training loss of the expert learning to replicate the aggregated mean fields generated by the FP-outputted policies (Algorithm \ref{algo:il_agg_meanfield}), and the training losses for the vanilla and adaptive parametric policies trained via Algorithm \ref{algo:interactive_imitation}.

% \begin{figure}
%     \centering
%     \includegraphics[width=0.9\linewidth]{Figures/BeachBar/beachbar_eta=0.3_FULL.pdf}
%     \caption{ Performance metrics for 5 different runs, with $\alpha = 1$.}
%     \label{fig:beachbar_alpha=1}
% \end{figure}

% \begin{figure}
%     \centering
%     \includegraphics[width=0.9\linewidth]{Figures/BeachBar/beachbar__alpha=1_FULL.pdf}
%     \caption{ Performance metrics for 5 different runs, with $\alpha = 1$.}
%     \label{fig:beachbar_alpha=1}
% \end{figure}
\begin{figure}[htpb]
    \centering
    % Left plot: Variation with eta
    \begin{subfigure}{0.49\linewidth}
        \centering
        \includegraphics[width=\linewidth]{Figures/BeachBar/beachbar_eta=0.3_FULL.pdf}
        \caption{Fixed $\eta = 0.3$.}
        \label{fig:beachbar_sub_eta}
    \end{subfigure}
    \hfill
    % Right plot: Variation with alpha
    \begin{subfigure}{0.49\linewidth}
        \centering
        \includegraphics[width=\linewidth]{Figures/BeachBar/beachbar__alpha=1_FULL.pdf}
        \caption{Fixed $\alpha = 1$.}
        \label{fig:beachbar_sub_alpha}
    \end{subfigure}

    \caption{Performance metrics for 5 runs in the Beach Bar environment. Left: variation across $\alpha$ with fixed $\eta$; Right: variation across $\eta$ with fixed $\alpha$.}
    \label{fig:beachbar_comparison_side_by_side}
\end{figure}

\begin{figure}[htbp]
    \centering
    \includegraphics[width=0.9\linewidth]{Figures/BeachBar/beachbar_GRID_FULL.pdf}
    \caption{Beach Bar: Metrics Averages}
    \label{fig:beachbar_grid_mean}
\end{figure}

\begin{figure}[htbp]
    \centering
    \includegraphics[width=0.9\linewidth]{Figures/BeachBar/beachbar_GRID_STD_FULL.pdf}
    \caption{Beach Bar: Metrics STD.}
    \label{fig:beachbar_grid_std}
\end{figure}


\begin{figure}[htbp]
    \centering
    \includegraphics[width=0.9\linewidth]{Figures/BeachBar/beachbar_losses_part_3.pdf}
    \caption{Beach Bar: Fictitious Play Exploitability Convergence, Expert Losses, Vanilla and Adaptive Losses ($\eta = 0.3)$.}
    \label{fig:beachbar_losses}
\end{figure}

\clearpage

\section{Additional Details to the Night Clubs Example}
\label{appendix:nightclubs}
\subsection{Hyperparameters}

All the experiments have been ran with the parameters presented in Table~\ref{tab:nightclubs_hyperparameters}
\begin{table}[htbp]
\centering
\caption{Summary of Hyperparameters for the Night Club Environment}
\label{tab:nightclubs_hyperparameters}
\begin{tabular}{@{}lll@{}}
\toprule
\textbf{Category} & \textbf{Parameter} & \textbf{Value} \\ \midrule
\textbf{Neural Architecture} & Model Type & MLP \\
& Hidden Layers & $2$ \\
& Layer Width & $64$ \\
& Activation Function & ReLU \\
& Optimizer & Adam \\ \addlinespace
\textbf{Fictitious Play} & Number of (FP) iterations & $40$ \\ \addlinespace
\textbf{Best Response Training} & Number of iterations & $1\,000$ \\
& Learning Rate & $10^{-4}$ \\
& Batch Size  & $500$ \\ \addlinespace
\textbf{Expert Imitation} & Number of iterations & $20\,000$ \\
& Learning Rate & $5\times 10^{-4}$ \\
& Scheduler & Cosine Annealing \\
& Batch Size & $500$ \\ \addlinespace
\textbf{Vanilla and Adaptive Learning} & Number of iterations & $20\,000$ \\
& Learning Rate & $10^{-4}$ \\
& Batch Size & $50$ \\ 
& Number of agents in the Simulation & $1000$ \\ 
\bottomrule
\end{tabular}
\end{table}

\subsection{Full Sensitivity Analysis and Training Losses}
Figure \ref{fig:nightclubs_grid_mean} represents the average over $5$ runs of the metrics, as functions of the varying parameters $\eta$ and $\alpha$. Figure \ref{fig:nightclubs_grid_std} shows the standard deviation. For the row corresponding to the proxy, the advantage corresponds to the adaptive proxy minus the vanilla proxy. For the rows corresponding to the reward vs expert, it corresponds to the reward of the adaptive policy minus the one of the vanilla policy. For the rows corresponding to the exploitability, it corresponds to the exploitability of the vanilla minus the one of the adaptive policy.

Figure \ref{fig:nightclubs_losses} illustrates the convergence of Fictitious Play (FP) with common noise (Algorithm \ref{algo:fictitious}), the training loss of the expert learning to replicate the aggregated mean fields generated by the FP-outputted policies (Algorithm \ref{algo:il_agg_meanfield}), and the training losses for the vanilla and adaptive parametric policies trained via Algorithm \ref{algo:interactive_imitation}.

% \begin{figure}[htbp]
%     \centering
%     \includegraphics[width=0.9\linewidth]{Figures/NightClubs/nightclub__alpha=0.1_FULL.pdf}
%     \caption{ Performance metrics for 5 different runs, with $\alpha = 0.1$.}
%     \label{fig:nightclubs_alpha=0.1}
% \end{figure}

\begin{figure}[htpb]
    \centering
    % Left plot: Variation with eta
    \begin{subfigure}{0.49\linewidth}
        \centering
        \includegraphics[width=\linewidth]{Figures/NightClubs/nightclub__alpha=1.05_FULL.pdf}
        \caption{Fixed $\alpha = 1.05$.}
        \label{fig:nightclubs_sub_alpha=1.05}
    \end{subfigure}
    \hfill
    % Right plot: Variation with alpha
    \begin{subfigure}{0.49\linewidth}
        \centering
        \includegraphics[width=\linewidth]{Figures/NightClubs/nightclub__alpha=0.1_FULL.pdf}
        \caption{Fixed $\alpha = 0.1$.}
        \label{fig:nightclubs_sub_alpha=0.1}
    \end{subfigure}

    \caption{Performance metrics for 5 different runs in the Night Clubs example, with $\alpha$ fixed and $\eta$ varying.}
    \label{fig:nightclubs_comparison_side_by_side}
\end{figure}

\begin{figure}[htbp]
    \centering
    \includegraphics[width=0.9\linewidth]{Figures/NightClubs/nightclub_GRID_FULL.pdf}
    \caption{Night Clubs: Metrics Averages}
    \label{fig:nightclubs_grid_mean}
\end{figure}

\begin{figure}[htbp]
    \centering
    \includegraphics[width=0.9\linewidth]{Figures/NightClubs/nightclub_GRID_STD_FULL.pdf}
    \caption{Night Clubs: Metrics STD.}
    \label{fig:nightclubs_grid_std}
\end{figure}


\begin{figure}[htbp]
    \centering
    \includegraphics[width=0.9\linewidth]{Figures/NightClubs/nightclub_losses_part_1.pdf}
    \caption{Night Clubs: Fictitious Play Exploitability Convergence, Expert Losses, Vanilla and Adaptive Losses ($\alpha = 0.1)$.}
    \label{fig:nightclubs_losses}
\end{figure}


% \clearpage 

% \color{blue}

% BELOW: ATTEMPT, NO NEED TO READ FOR NOW

% \section{Formal Reconstruction Framework}

% We address the problem of reconstructing the generating policy of a MFG from observed trajectories. We focus on the ``inverse problem'' at a specific time step. Let $\mu \in \Delta_{\cX}$ be a fixed population measure (with $\mu(x) > 0$ for all $x \in \cX$) and let $\pi^* \in \Pi$ be the unknown true policy employed by the population.

% We assume we observe a dataset $\mathcal{D}_K = \{ (\varepsilon^{(k)}, \muNext^{(k)}) \}_{k=1}^K$ consisting of $K$ independent realizations of the common noise $\varepsilon \sim \nu^0$ and the resulting next-step measure $\muNext$.

% \subsection{Vectorization and Dynamics}
% We represent the local policy $\pi(\cdot | \cdot, \mu)$ as a vector $\bz \in \RR^{|\cX||\cA|}$ indexed by state-action pairs $(x,a)$, such that $z_{(x,a)} = \pi(a|x, \mu)$. The set of valid policies is the compact convex polytope $\mathcal{Z}$:
% \begin{equation}
%     \mathcal{Z} = \left\{ \bz \in \RR^{|\cX||\cA|} \;\middle|\; \bz \succeq 0, \quad \sum_{a \in \cA} z_{(x,a)} = 1, \forall x \in \cX \right\}.
% \end{equation}

% The dynamics of the mean field are linear with respect to the policy vector $\bz$, though stochastic due to the common noise. We define the stochastic transition matrix $\bM(\varepsilon) \in \RR^{|\cX| \times |\cX||\cA|}$ for a realization $\varepsilon$ as:
% \begin{equation}
%     [\bM(\varepsilon)]_{y, (x,a)} \coloneqq \mu(x) P(y \mid x, a, \mu, \varepsilon).
% \end{equation}
% The observation model is then given by:
% \begin{equation}
%     \muNext^{(k)} = \bM(\varepsilon^{(k)}) \bz^*, \quad \text{for } k=1, \dots, K.
% \end{equation}

% \subsection{Estimator Definition}
% We propose a \textit{Constrained Least Squares Estimator} (CLSE). We seek the policy that minimizes the mean squared prediction error over the observed common noise realizations.
% Define the empirical risk function $L_K: \mathcal{Z} \to \RR_+$:
% \begin{equation}
%     L_K(\bz) \coloneqq \frac{1}{K} \sum_{k=1}^K \big\| \bM(\varepsilon^{(k)}) \bz - \muNext^{(k)} \big\|_2^2.
% \end{equation}
% The estimator $\hZ_K$ is defined as:
% \begin{equation}
%     \label{eq:estimator}
%     \hZ_K \in \argmin_{\bz \in \mathcal{Z}} L_K(\bz).
% \end{equation}
% Since $\mathcal{Z}$ is compact and $L_K$ is continuous (and convex), a solution always exists.

% \subsection{Convergence Analysis}

% We analyze the asymptotic behavior of $\hZ_K$ as $K \to \infty$. Let $L(\bz)$ denote the expected (population) risk:
% \begin{equation}
%     L(\bz) \coloneqq \EE_{\varepsilon \sim \nu^0} \left[ \big\| \bM(\varepsilon) (\bz - \bz^*) \big\|_2^2 \right] = (\bz - \bz^*)^\top \bQ (\bz - \bz^*),
% \end{equation}
% where $\bQ \coloneqq \EE [ \bM(\varepsilon)^\top \bM(\varepsilon) ]$ is the expected Gram matrix of the transition dynamics.

% \begin{assumption}[Identifiability]
%     \label{ass:identifiability}
%     The matrix $\bQ$ is positive definite on the tangent space of the simplex constraints. That is, for any deviation vector $\bv \in \RR^{|\cX||\cA|} \setminus \{0\}$ satisfying $\sum_{a} v_{(x,a)} = 0$ for all $x$, we have:
%     \begin{equation}
%         \bv^\top \bQ \bv > 0.
%     \end{equation}
% \end{assumption}

% \textit{Remark:} This assumption implies that the common noise sufficiently randomizes the transitions such that no two distinct policies produce the same expected transition dynamics almost surely.

% \begin{theorem}[Strong Consistency]
%     Let the Identifiability Assumption \ref{ass:identifiability} hold. Then the sequence of estimators $\hZ_K$ converges almost surely to the true policy $\bz^*$ as $K \to \infty$:
%     \begin{equation}
%         \hZ_K \xrightarrow{a.s.} \bz^*.
%     \end{equation}
% \end{theorem}

% \begin{proof}
%     The proof follows standard consistency arguments for M-estimators on compact sets.
    
%     1. \textbf{Uniform Convergence:} The parameter space $\mathcal{Z}$ is compact. The loss function $((\varepsilon, \by), \bz) \mapsto \| \bM(\varepsilon)\bz - \by \|^2$ is continuous and bounded on $\mathcal{Z}$ (assuming bounded dynamics). By the Uniform Law of Large Numbers (specifically, applying the uniform convergence of sample means over compact sets), we have:
%     \begin{equation}
%         \sup_{\bz \in \mathcal{Z}} | L_K(\bz) - L(\bz) | \xrightarrow{a.s.} 0.
%     \end{equation}
    
%     2. \textbf{Unique Minimizer:} clearly $L(\bz^*) = 0$. Under Assumption \ref{ass:identifiability}, $L(\bz)$ is strictly convex restricted to valid directions in $\mathcal{Z}$. Thus, $\bz^*$ is the unique global minimizer of $L$ on $\mathcal{Z}$.
    
%     3. \textbf{Convergence of Argmin:} Since $L_K$ converges uniformly to $L$ and $\bz^*$ is the unique well-separated minimum of $L$, it follows that the sequence of minimizers $\hZ_K$ converges to $\bz^*$. Formally, for any $\delta > 0$, let $\gamma = \inf \{ L(\bz) : \bz \in \mathcal{Z}, \| \bz - \bz^* \| \ge \delta \}$. By strict convexity, $\gamma > 0$. Uniform convergence implies that for large enough $K$, $|L_K(\bz) - L(\bz)| < \gamma/3$ for all $\bz$. This ensures that the minimum of $L_K$ cannot lie outside the $\delta$-ball of $\bz^*$. Thus $\hZ_K \to \bz^*$ almost surely.
% \end{proof}




\end{document}

% This document was modified from the file originally made available by
% Pat Langley and Andrea Danyluk for ICML-2K. This version was created
% by Iain Murray in 2018, and modified by Alexandre Bouchard in
% 2019 and 2021 and by Csaba Szepesvari, Gang Niu and Sivan Sabato in 2022.
% Modified again in 2023 and 2024 by Sivan Sabato and Jonathan Scarlett.
% Previous contributors include Dan Roy, Lise Getoor and Tobias
% Scheffer, which was slightly modified from the 2010 version by
% Thorsten Joachims & Johannes Fuernkranz, slightly modified from the
% 2009 version by Kiri Wagstaff and Sam Roweis's 2008 version, which is
% slightly modified from Prasad Tadepalli's 2007 version which is a
% lightly changed version of the previous year's version by Andrew
% Moore, which was in turn edited from those of Kristian Kersting and
% Codrina Lauth. Alex Smola contributed to the algorithmic style files.
